% archon:covers AlgebraicJacobian/RiemannRoch/OcOfD.lean
% archon:covers AlgebraicJacobian/RiemannRoch/SmoothStalkDVR.lean

\chapter{The invertible sheaf \(\mathcal O_C(D)\) of a Weil divisor (RR.2\(_\ast\))}
\label{chap:RiemannRoch_OcOfD}

% SOURCE: Hartshorne, \emph{Algebraic Geometry}, II~\S6, pp.~144--145
% (definition of \(\mathscr L(D)\); Proposition 6.13; Remark 6.17.1;
% Proposition 6.18) and IV.1, p.~296 (the inductive-step short exact
% sequence). Read directly from
% references/hartshorne-algebraic-geometry.pdf, PDF pages 161--162 and
% 313, by rendering the scanned page images and transcribing the
% statements character-by-character (the PDF has no text layer). Stacks
% pointers (tag-level, not verbatim-quoted): 02RW (Weil divisors), 0AUW
% (the sheaf \(\mathcal O_X(D)\) from a Weil divisor), 0BE3 (degree-zero of
% a principal divisor on a complete nonsingular curve).

\section{Setup and motivation}

This chapter is a satellite of \cref{chap:RiemannRoch_RRFormula};
together with \cref{chap:RiemannRoch_WeilDivisor} and
\cref{chap:RiemannRoch_OCofP}, it supplies the sheaf-theoretic
infrastructure for the Riemann--Roch bridge: a smooth proper geometrically
irreducible curve \(C / \bar k\) with \(g(C) = 0\) is isomorphic to
\(\mathbb P^1_{\bar k}\) (\cref{thm:genus_zero_curve_iso_p1}).

The genus-zero rigidity argument relies on the invertible-sheaf functor
\[
  \mathcal O_C(-) \;\colon\; \mathrm{Div}(C) \;\longrightarrow\;
  \mathbf{Sh}(C, \mathbf{Mod}_{\bar k})
\]
sending a Weil divisor \(D = \sum_i n_i [P_i]\) to its associated invertible
sheaf \(\mathcal O_C(D)\) on \(C\). This chapter constructs that functor and
its three immediate corollaries: the structure-sheaf identification
\(\mathcal O_C(0) = \mathcal O_C\) (\cref{lem:sheafOf_zero}), the
closed-point comparison \(\mathcal O_C([P]) \cong
\texttt{lineBundleAtClosedPoint}_P\) (\cref{def:lineBundleAtClosedPoint},
\cref{lem:sheafOf_singlePoint}), and the inductive-step short exact
sequence \(0 \to \mathcal O_C(D) \to \mathcal O_C(D+[Y]) \to k(Y) \to 0\)
(\cref{lem:sheafOf_ses_single_add}).

\paragraph{Relation to the Riemann--Roch induction.}
The \(\chi\)-identity proof of
\cref{thm:euler_char_eq_deg_plus_one_minus_genus} in
\cref{chap:RiemannRoch_RRFormula} is an induction on the Weil divisor
\(D\) via the \(D \rightsquigarrow D + [Y]\) step of Hartshorne~IV.1.3.
That step requires the short exact sequence
\[
  0 \to \mathcal O_C(D) \to \mathcal O_C(D + [Y]) \to k(Y)^{[\text{at } Y]}
  \to 0
\]
to be exact in \(\mathbf{Coh}(C)\); this is a property of the
\emph{body} of the functor \(\mathcal O_C(-)\), not merely of its
signature, and is the content of \cref{lem:sheafOf_ses_single_add}.

\paragraph{Relation to the line bundle at a closed point.}
The chapter \cref{chap:RiemannRoch_OCofP} constructs the line bundle
\(\mathcal O_C([P])\) associated to a \emph{closed point} \(P \in C\)
directly, bypassing the Weil-divisor packaging.
The \cref{lem:sheafOf_singlePoint} below makes this construction
mathematically coherent with the general functor by exhibiting
\(\mathcal O_C(D)\) at \(D = [P]\) as canonically \emph{isomorphic} to
\cref{def:lineBundleAtClosedPoint}, via an isomorphism that is the identity
on the common carrier \(K(C)\). This comparison enables the
downstream use of \(\dim_{\bar k} H^0(C, \mathcal O_C(P)) = 2\)
via the \(\chi\)-identity.

\paragraph{Relation to the rational-curve isomorphism.}
The chapter \cref{chap:RiemannRoch_RationalCurveIso} translates
Riemann--Roch outputs into the function-field/degree language of \(C\);
this translation references \(\mathcal O_C(D)\) for general \(D\) (not
just \(D = [P]\)). The function-field-degree bridge reads off \(\deg(D)\)
from the associated invertible sheaf, and the Weil-divisor / line-bundle
correspondence (Hartshorne's Proposition~6.15) is the vehicle for that
reading, pinned here via \cref{lem:sheafOf_singlePoint}.

\paragraph{Standing hypotheses.} Throughout this chapter, \(\bar k\) is an
algebraically closed field and \(C\) is a smooth proper geometrically
irreducible curve over \(\bar k\). The Weil divisor group
\(\mathrm{Div}(C)\) is the one defined in \cref{def:codim1_cycles}; the
prime-divisor notion is \cref{def:prime_divisor}; the order
\(\ord_Q(f)\) at a prime divisor \(Q\) is \cref{def:order_at_point}; the
principal divisor \(\div(f)\) of a non-zero rational function is
\cref{def:principal_divisor}; the closed-point divisor \([P]\) is
\cref{def:divisor_closed_point}. The line bundle at a closed point is
\cref{def:lineBundleAtClosedPoint}, and \cref{lem:sheafOf_singlePoint}
makes the comparison
\(\mathcal O_C(D)|_{D = [P]} \cong \texttt{lineBundleAtClosedPoint}_P\)
explicit.

\paragraph{Construction strategy.}
The functor \(\mathcal O_C(D)\) is defined following Hartshorne~II~\S6 as
the sub-\(\mathcal O_C\)-module of the constant sheaf
\(\mathcal K_C \cong K(C)\) generated by the appropriate local data. Both
\(\mathcal O_C(D)|_{D=[P]}\) and \cref{def:lineBundleAtClosedPoint}
use Hartshorne's subsheaf-of-\(\mathcal K_C\) recipe with the same section
carriers, so the comparison of \cref{lem:sheafOf_singlePoint} is the
identity-on-\(K(C)\) isomorphism induced by the section-wise equality of
those carriers.

\section{The invertible sheaf \(\mathcal O_C(D)\)}

\begin{definition}\leanok
[Invertible sheaf \(\mathcal O_C(D)\) of a Weil divisor]
  \label{def:sheafOf}
  \lean{AlgebraicGeometry.Scheme.WeilDivisor.sheafOf}
  \uses{def:codim1_cycles, def:prime_divisor, def:order_at_point,
        def:Scheme_toModuleKPresheaf, thm:Scheme_toModuleKPresheaf_isSheaf,
        def:sheafOf_carrierPresheaf, lem:sheafOf_carrierPresheaf_isSheaf,
        thm:smooth_isRegularInCodimOne}

  % SOURCE: [Hartshorne], II.6, p.~144 (definition of the invertible
  % sheaf \(\mathscr L(D)\) associated to a Cartier divisor on a scheme),
  % and Remark~6.17.1, Proposition~6.18 + Proposition~6.15, p.~145
  % (Cartier \(\leftrightarrow\) Weil correspondence on an integral
  % locally factorial scheme; identification of \(\mathscr L(D)\) as a
  % subsheaf of the constant sheaf at the function field \(K\)). Read from
  % references/hartshorne-algebraic-geometry.pdf, PDF pages 161--162,
  % transcribed character-by-character from rendered scanned page
  % images (the PDF has no text layer).
  % SOURCE QUOTE: "Definition. Let \(D\) be a Cartier divisor on a scheme
  % \(X\), represented by \(\{(U_i, f_i)\}\) as above. We define a subsheaf
  % \(\mathscr L(D)\) of the sheaf of total quotient rings \(\mathscr K\) by
  % taking \(\mathscr L(D)\) to be the sub-\(\mathscr O_X\)-module of
  % \(\mathscr K\) generated by \(f_i^{-1}\) on \(U_i\). This is well-defined,
  % since \(f_i/f_j\) is invertible on \(U_i \cap U_j\), so \(f_i^{-1}\) and
  % \(f_j^{-1}\) generate the same \(\mathscr O_X\)-module. We call
  % \(\mathscr L(D)\) the \emph{sheaf associated to \(D\)}."
  % SOURCE QUOTE: "Proposition 6.15. \emph{If \(X\) is an integral
  % scheme, the homomorphism \(\mathrm{CaCl}\, X \to \mathrm{Pic}\, X\) of
  % (6.14) is an isomorphism.} Proof. We have only to show that every
  % invertible sheaf is isomorphic to a subsheaf of \(\mathscr K\), which
  % in this case is the constant sheaf \(K\), where \(K\) is the function
  % field of \(X\)."
  % SOURCE QUOTE: "Remark 6.17.1. Clearly this gives a 1-1
  % correspondence between effective Cartier divisors on \(X\) and
  % \emph{locally principal} closed subschemes \(Y\), i.e., subschemes
  % whose sheaf of ideals is locally generated by a single element.
  % Note also that if \(X\) is an integral separated noetherian locally
  % factorial scheme, so that the Cartier divisors correspond to Weil
  % divisors by (6.11), then the effective Cartier divisors correspond
  % exactly to the effective Weil divisors."
  \textit{Source: Hartshorne, II.6, p.~144 (the definition of
  \(\mathscr L(D)\)), Proposition~6.15 (the constant-sheaf-at-\(K\)
  identification on an integral scheme), and Remark~6.17.1 (Weil
  divisors and Cartier divisors agree on an integral separated
  noetherian locally factorial scheme).}
  Let \(C\) be a smooth proper geometrically irreducible curve over
  \(\bar k\) and let \(D = \sum_{Q} n_Q [Q] \in \mathrm{Div}(C)\) be a Weil divisor
  on \(C\), with the formal sum indexed by prime divisors \(Q\) of \(C\) and
  supported on a finite set \(\{Q : n_Q \neq 0\}\)
  (\cref{def:codim1_cycles}). The \emph{invertible sheaf
  \(\mathcal O_C(D)\) associated to \(D\)}, also written
  \(\mathcal O_C(D) := \mathscr L(D)\) in Hartshorne's notation, is the
  sub-\(\mathcal O_C\)-module of the constant sheaf
  \(\mathcal K_C \cong K(C)\) (Proposition~6.15: on an integral scheme the
  sheaf of total quotient rings is the constant sheaf at the function
  field) defined section-wise by
  \[
    \Gamma\bigl(U, \mathcal O_C(D)\bigr) \;:=\;
    \{\, 0 \,\} \;\cup\;
    \bigl\{\, f \in K(C)^{\times} \;\big|\;
      \ord_Q(f) \;\geq\; -\,n_Q \text{ for every prime divisor }
      Q \text{ of } C \text{ with } Q \in U \bigr\}
  \]
  for each open \(U \subseteq C\). Restriction along \(V \hookrightarrow U\)
  is the identity on the underlying carrier \(K(C)\), which preserves
  membership (a smaller open imposes \emph{fewer} order constraints, so
  the constraint set on \(U\) embeds into the constraint set on \(V\)). The
  \(\bar k\)-module structure on
  \(\Gamma(U, \mathcal O_C(D))\) is restricted from the
  \(\bar k\)-module structure on \(K(C)\) (a \(\bar k\)-scalar multiple of a
  rational function has the same order at every prime divisor).

  \paragraph{Equivalence with the Cartier-divisor description.} The
  formal sum \(D = \sum_Q n_Q [Q]\) packages, on a curve where prime
  divisors are exactly closed points and the local rings at prime
  divisors are DVRs (smoothness \(+\) dimension~\(1\)), the Cartier-divisor
  data \(\{(U_i, f_i)\}\) of Hartshorne's definition: locally on a small
  open \(U_i\) containing exactly one prime divisor \(Q\) in the support of
  \(D\), pick a uniformiser \(\pi_Q \in K(C)\) at \(Q\) and set
  \(f_i := \pi_Q^{n_Q}\); on opens disjoint from the support of \(D\), set
  \(f_i := 1\). The transitions on overlaps \(U_i \cap U_j\) are units in
  \(\mathcal O_{C, Q}^{\times}\) for each prime \(Q\) in the overlap (any
  two uniformisers differ by a unit, by the DVR structure), so the
  \(\{(U_i, f_i)\}\) define a Cartier divisor whose Weil-divisor image
  under Hartshorne's (6.11) is the original \(D\). The
  ``\(f_i^{-1}\)-generated'' subsheaf of \(\mathcal K_C\) from Hartshorne's
  definition then coincides with the section-wise description above:
  the local generator \(f_i^{-1} = \pi_Q^{-n_Q}\) has order \(-n_Q\) at \(Q\),
  so the local sections are exactly \(f \in K(C)\) with
  \(\ord_Q(f \cdot \pi_Q^{n_Q}) \geq 0\), i.e.\ \(\ord_Q(f) \geq -n_Q\), on
  \(U_i\).

  Hartshorne's Proposition~6.13(a) (see
  \cref{sec:sheaf-property-correctness})
  then guarantees that the resulting subsheaf $\mathcal O_C(D) \subset
  \mathcal K_C$ is an invertible sheaf on $C$, independently of the
  choice of uniformisers \(\pi_Q\) (any two choices differ by a unit, so
  they generate the same \(\mathcal O_C\)-submodule).

  \paragraph{Construction overview.}
  The closed-point specialisation \(\mathcal O_C([P])\) is
  \cref{def:lineBundleAtClosedPoint} of \cref{chap:RiemannRoch_OCofP};
  their agreement at \(D = [P]\) is \cref{lem:sheafOf_singlePoint} below.
  The construction proceeds presheaf-by-hand via per-open
  \(\bar k\)-submodules of the function field, with the sheaf property
  reduced to the locality of the order conditions at the stalks of the
  prime divisors. The substrate helpers
  (\cref{def:sheafOf_carrierSet,lem:sheafOf_carrierSet_mono,%
  def:sheafOf_carrierSubmodule,def:sheafOf_trivAtBot,%
  def:sheafOf_carrierSubmoduleSheaf,lem:sheafOf_carrierSubmoduleSheaf_le,%
  def:sheafOf_carrierTypeSubfunctor,def:sheafOf_carrierPresheaf,%
  lem:sheafOf_carrierPresheaf_isSheaf}) are assembled in the
  ``Carrier substrate'' section below.

  \paragraph{Substrate obligation.} The general (\(D \neq 0\)) branch of
  the construction relies on the hypothesis that \(C\) is regular in
  codimension one, i.e.\ that every local ring \(\mathcal O_{C, Q}\) at a
  prime divisor \(Q\) is a discrete valuation ring. This is the smoothness
  \(\Rightarrow\) DVR implication at a \(\bar k\)-rational codimension-one
  point, obtained through the one-dimensional cotangent space and
  established in \cref{thm:smooth_isRegularInCodimOne}.
  Every carrier helper below is established outright once this DVR
  substrate is in scope.
\end{definition}

\section{The codimension-one regularity substrate}

The general (\(D \neq 0\)) branch of \cref{def:sheafOf} forms its order
conditions inside the local rings \(\mathcal O_{C, Q}\) at the prime
divisors \(Q\) of \(C\), and reads off \(\ord_Q\) as the valuation of a
discrete valuation ring; this is sound precisely when \(C\) is regular in
codimension one (\cref{def:isRegularInCodimensionOne}). The instance below
establishes this from the curve's smoothness hypothesis,
generalising the \(\mathbb P^1\)-only
\cref{thm:projectiveLineBar_isRegularInCodimOne} of
\cref{chap:RiemannRoch_WeilDivisor} to an abstract smooth curve \(C\).

The substrate is established via the conormal isomorphism. At a codimension-one point \(x \in C\) the stalk \(\mathcal O_{C, x}\)
is a smooth local \(\bar k\)-algebra: it is the localisation at a height-one
prime of a relative-dimension-one standard-smooth affine chart of \(C\), hence
essentially of finite type and formally smooth over \(\bar k\). Since \(x\) has
coheight one and the curve has Krull dimension at most one
(\cref{thm:krullDim_curve_le_one}), the point \(x\) is closed, so its prime in
the standard-smooth chart is maximal; the structure map onto the residue field
is then surjective (\cref{lem:algebraMap_residueField_surjective}), exhibiting
\(\kappa(x)\) as a finite-type \(\bar k\)-algebra that is a field, hence
\(\kappa(x) = \bar k\) is \(\bar k\)-rational. The conormal isomorphism of
\cref{chap:RiemannRoch_SmoothRegular},
\[
  \mathfrak m_x / \mathfrak m_x^2 \;\cong\;
  \Omega_{C/\bar k, x} \otimes_{\mathcal O_{C, x}} \kappa(x)
  \qquad (\text{Hartshorne~II.8.7}),
\]
then identifies the residue-field dimension of the cotangent space with the
\(\mathcal O_{C, x}\)-rank of the differentials:
\(\operatorname{finrank}_{\kappa(x)} \mathfrak m_x / \mathfrak m_x^2
= \operatorname{finrank}_{\mathcal O_{C, x}} \Omega_{C/\bar k, x}\)
(\cref{lem:finrank_cotangentSpace_eq_rank_kaehler}). The differentials
\(\Omega_{C/\bar k}\) are locally free of rank one on the smooth curve
(\cref{lem:smooth_kaehler_locallyFree_rank_one}), so their stalk has rank one,
giving
\(\operatorname{finrank}_{\kappa(x)} \mathfrak m_x / \mathfrak m_x^2 = 1\)
(\cref{lem:cotangentSpace_finrank_eq_one_of_smooth}).
\cref{lem:isDVR_stalk_of_finrank_cotangentSpace} (via
\cref{lem:finrank_cotangentSpace_eq_one_iff_isDVR}) then yields
\(\texttt{IsDiscreteValuationRing}\,\mathcal O_{C, x}\)
(\cref{lem:stalk_isDVR_of_smooth}), and
\cref{thm:smooth_isRegularInCodimOne} applies this at each
codimension-one point. The one-dimensional cotangent dimension is obtained
through the conormal identification; the only use of dimension data is the
global curve bound \(\operatorname{krullDim} C.\mathrm{left} \le 1\)
(\cref{thm:krullDim_curve_le_one}), which serves solely to certify that each
codimension-one point is closed (and hence \(\bar k\)-rational) and supplies the
Krull-dimension-bound hypothesis of \cref{lem:cotangentSpace_finrank_eq_one_of_smooth}.
This is the abstract-curve generalisation of
\cref{thm:projectiveLineBar_isRegularInCodimOne}.

\subsection{The abstract ring-bridge (algebra-core realisation)}

The conormal route above is realised at the level of commutative algebra by
three lemmas about a \emph{localisation} \(T\) of a relative-dimension-one
standard-smooth \(k\)-algebra \(S\). At a codimension-one point \(x \in C\) the
stalk \(\mathcal O_{C, x}\) is exactly such a \(T\): the localisation
\(T = S_{\mathfrak q}\) of the coordinate ring \(S\) of a standard-smooth affine
chart of \(C\) at the height-one prime \(\mathfrak q\) corresponding to \(x\).
These three lemmas form the pure commutative-algebra core consumed by the
scheme-level statements of this section.

\begin{lemma}\leanok
[Free rank-one differentials at a standard-smooth localisation]
  \label{lem:free_finrank_kaehler_localization}
  \lean{Algebra.free_and_finrank_kaehlerDifferential_localization_eq_one}
  Let \(k\) be a field, let \(S\) be a relative-dimension-one standard-smooth
  \(k\)-algebra (\(\texttt{IsStandardSmoothOfRelativeDimension}\ 1\ k\ S\)), and
  let \(T\) be a localisation of \(S\) at a submonoid \(M \subseteq S\). Then the
  module of Kähler differentials \(\Omega_{T/k}\) is a free \(T\)-module of rank
  one:
  \[
    \texttt{Module.Free}\ T\ \Omega_{T/k}
    \qquad\text{and}\qquad
    \operatorname{finrank}_T \Omega_{T/k} = 1.
  \]
  This is the algebra realisation consumed in the stalk-level half of
  \cref{lem:smooth_kaehler_locallyFree_rank_one}.
\end{lemma}

\begin{proof}

\leanok
  The standard-smooth \(k\)-algebra \(S\) has free Kähler differentials
  \(\Omega_{S/k}\), of rank equal to the relative dimension, here one. Kähler
  differentials commute with localisation: the canonical map
  \(\Omega_{S/k} \to \Omega_{T/k}\) exhibits \(\Omega_{T/k}\) as the
  localisation of \(\Omega_{S/k}\) at \(M\).
  A localisation of a
  free rank-one module is free of rank one, whence \(\Omega_{T/k}\) is
  \(T\)-free with \(\operatorname{finrank}_T \Omega_{T/k} = 1\).
\end{proof}

\begin{lemma}

\leanok
[One-dimensional cotangent space at a standard-smooth localisation]
  \label{lem:finrank_cotangent_localization_eq_one}
  \lean{Algebra.finrank_cotangentSpace_localization_eq_one_of_isStandardSmooth_dim_one}
  \uses{lem:free_finrank_kaehler_localization,
        lem:finrank_cotangentSpace_eq_rank_kaehler}
  With \(k, S, M, T\) as in \cref{lem:free_finrank_kaehler_localization}, suppose
  in addition that \(T\) is a local ring (\(\texttt{IsLocalRing}\ T\)) whose
  residue field \(\kappa(T)\) is \(k\)-rational --- with \(\Omega_{\kappa(T)/k} = 0\)
  and \(H^1_{\mathrm{cot}}(k, \kappa(T)) = 0\) (both hold when
  \(\kappa(T) = k\)). Then the cotangent space of \(T\) is one-dimensional over
  its residue field:
  \[
    \operatorname{finrank}_{\kappa(T)}
      \bigl(\texttt{IsLocalRing.CotangentSpace}\ T\bigr) = 1.
  \]
\end{lemma}

\begin{proof}
  \uses{lem:free_finrank_kaehler_localization,
        lem:finrank_cotangentSpace_eq_rank_kaehler}
        \leanok
  By \cref{lem:free_finrank_kaehler_localization} the differentials
  \(\Omega_{T/k}\) are \(T\)-free of rank one. The conormal corollary
  \cref{lem:finrank_cotangentSpace_eq_rank_kaehler},
  whose \(k\)-rationality hypotheses are supplied by the two residue-field
  vanishing assumptions, equates the residue-field dimension of the cotangent
  space with the \(T\)-rank of the differentials,
  \[
    \operatorname{finrank}_{\kappa(T)}
      \bigl(\texttt{CotangentSpace}\ T\bigr)
    = \operatorname{finrank}_T \Omega_{T/k} = 1.
  \]
  The two residue-field vanishing assumptions are exactly the \(k\)-rationality
  hypotheses this lemma carries; in the scheme-level application they are
  supplied at a codimension-one point of a smooth curve, where the residue
  field is the base field \(\bar k\) (see the proof of
  \cref{lem:cotangentSpace_finrank_eq_one_of_smooth}).
\end{proof}

\begin{lemma}\leanok
[Discrete valuation ring at a standard-smooth localisation]
  \label{lem:isDVR_localization_standardSmooth}
  \lean{Algebra.isDiscreteValuationRing_localization_of_isStandardSmooth_dim_one}
  \uses{lem:finrank_cotangent_localization_eq_one}
  With \(k, S, M, T\) as in \cref{lem:finrank_cotangent_localization_eq_one},
  suppose moreover that \(T\) is a Noetherian local integral domain
  (\(\texttt{IsLocalRing}\ T\), \(\texttt{IsNoetherianRing}\ T\),
  \(\texttt{IsDomain}\ T\)) with \(k\)-rational residue field. Then \(T\) is a
  discrete valuation ring (\(\texttt{IsDiscreteValuationRing}\ T\)).
\end{lemma}

\begin{proof}
  \uses{lem:finrank_cotangent_localization_eq_one}
  \leanok
  As a localisation of the finite-type \(k\)-algebra \(S\), the ring \(T\) is
  essentially of finite type over \(k\); and standard-smoothness of \(S\) makes
  \(S\) formally smooth over \(k\), which composes with the formal smoothness of
  the localisation map \(S \to T\) to make \(T\) formally smooth over \(k\). Thus
  \(T\) is a smooth local \(k\)-algebra of relative dimension one, whose
  cotangent space is one-dimensional by
  \cref{lem:finrank_cotangent_localization_eq_one}. The smooth-DVR criterion
  of \cref{chap:RiemannRoch_SmoothRegular} then yields
  \(\texttt{IsDiscreteValuationRing}\ T\).
\end{proof}

% --- Shared chart-extraction frontier -----------------------------------------
% The three scheme-level pins below
% (\cref{lem:smooth_kaehler_locallyFree_rank_one},
%  \cref{lem:cotangentSpace_finrank_eq_one_of_smooth},
%  \cref{lem:stalk_isDVR_of_smooth}) all need the SAME geometric input: a
% relative-dimension-one standard-smooth affine chart at the point $x$, with its
% stalk realised as a localisation over the base field $\bar k$. The next two
% lemmas isolate that input once; each pin then \uses{} them together with the
% pure-algebra core above.

\begin{lemma}

\leanok
[Standard-smooth affine chart over the base field at a point]
  \label{lem:exists_standardSmooth_chart_of_smooth}
  \lean{AlgebraicGeometry.exists_isStandardSmooth_chart_of_smooth}
  Let \(\bar k\) be a field and let \(C\) be a scheme over
  \(\operatorname{Spec} \bar k\) whose structure morphism is
  \(\texttt{SmoothOfRelativeDimension}\ 1\). Then for every point
  \(x \in C\) there is an affine open \(V \ni x\) with coordinate
  ring \(S = \Gamma(C, V)\) and a ring homomorphism
  \(\psi : \bar k \to S\) realising \(S\) as a base-field algebra whose
  structure map is standard-smooth of relative dimension one:
  \[
    \texttt{RingHom.IsStandardSmoothOfRelativeDimension}\ 1\ \psi .
  \]
\end{lemma}

\begin{proof}

\leanok
  The smoothness hypothesis provides, for the given \(x\), an affine open
  \(U \subseteq \operatorname{Spec} \bar k\) and an affine open
  \(V \ni x\) of \(C\) mapping into \(U\) under the structure morphism, together
  with a proof that the induced ring map
  \(\Gamma(\operatorname{Spec} \bar k, U) \to \Gamma(C, V)\)
  is standard-smooth of relative dimension one.

  Since \(\bar k\) is a field, \(\operatorname{Spec} \bar k\) is a single point,
  so the only affine open is the whole space \(U = \operatorname{Spec} \bar k\),
  and the canonical isomorphism \(\Gamma(\operatorname{Spec} \bar k) \cong \bar k\)
  identifies the source of the chart map with \(\bar k\). Transporting
  standard-smoothness across this isomorphism yields the base-field structure map
  \(\psi : \bar k \to S\) with
  \(\texttt{RingHom.IsStandardSmoothOfRelativeDimension}\ 1\ \psi\).
\end{proof}

\begin{lemma}

\leanok
[Stalk of a smooth curve as a base-field localisation of a standard-smooth chart]
  \label{lem:stalk_standardSmooth_localization_of_smooth}
  \lean{AlgebraicGeometry.isLocalization_stalk_standardSmooth_chart_of_smooth}
  \uses{lem:exists_standardSmooth_chart_of_smooth}
  With \(\bar k\), \(C\), \(x\), the affine chart \(V \ni x\),
  \(S = \Gamma(C, V)\), and \(\psi : \bar k \to S\) as in
  \cref{lem:exists_standardSmooth_chart_of_smooth}, let
  \(\mathfrak q \subseteq S\) be the prime ideal corresponding to \(x\) and let
  \(T = \mathcal O_{C, x}\) be the stalk. Then \(T\) is the localisation of
  \(S\) at \(\mathfrak q\), and the localisation and base-field structures are
  compatible: \(S\) and \(T\) carry \(\bar k\)-algebra structures forming a
  scalar tower \(\bar k \to S \to T\), the identification
  \(\texttt{IsLocalization}\ (S \setminus \mathfrak q)\ T\) holds, and \(S\) is
  standard-smooth of relative dimension one over \(\bar k\)
  (\(\texttt{Algebra.IsStandardSmoothOfRelativeDimension}\ 1\ \bar k\ S\)).
  This is precisely the algebraic data consumed by the pure-algebra core
  \cref{lem:free_finrank_kaehler_localization},
  \cref{lem:finrank_cotangent_localization_eq_one},
  \cref{lem:isDVR_localization_standardSmooth}.
\end{lemma}

\begin{proof}
  \uses{lem:exists_standardSmooth_chart_of_smooth}
  \leanok
  Because \(V\) is an affine open containing \(x\), the stalk
  \(T = \mathcal O_{C, x}\) is the localisation of its coordinate ring
  \(S = \Gamma(C, V)\) at the prime \(\mathfrak q\) corresponding to \(x\),
  giving an \(S\)-algebra structure on \(T\) and the localisation identification
  \(T \cong S_{\mathfrak q}\). Equipping \(S\) and \(T\) with their base-field
  \(\bar k\)-algebra structures via the chart map \(\psi : \bar k \to S\) from
  \cref{lem:exists_standardSmooth_chart_of_smooth} and the composite
  \(\bar k \to S \to T\) yields the scalar tower \(\bar k \to S \to T\); the
  standard-smoothness of the chart map \(\psi\) transports to the algebra-level
  instance \(\texttt{Algebra.IsStandardSmoothOfRelativeDimension}\ 1\ \bar k\ S\).
\end{proof}

\begin{lemma}[Smooth of relative dimension one \(\Rightarrow\) smooth]
  \label{lem:smoothOfRelativeDimension_smooth}
  \lean{AlgebraicGeometry.SmoothOfRelativeDimension.smooth}
  \mathlibok
  \textit{Provided by Mathlib
  (\texttt{Mathlib.AlgebraicGeometry.Morphisms.Smooth}).}
  If a morphism of schemes \(f : X \to Y\) is smooth of relative dimension
  \(n\) (\(\texttt{SmoothOfRelativeDimension}\ n\ f\)), then \(f\) is smooth
  (\(\texttt{Smooth}\ f\)). In particular the structure morphism
  \(C.\mathrm{hom} : C \to \operatorname{Spec} \bar k\) of a curve that is
  \(\texttt{SmoothOfRelativeDimension}\ 1\) is a smooth morphism.
\end{lemma}

\begin{lemma}
[Smooth curve \(\Rightarrow\) differentials locally free of rank one]
  \label{lem:smooth_kaehler_locallyFree_rank_one}
  \lean{AlgebraicGeometry.kaehlerDifferential_locallyFree_rank_one_of_smooth}
  \uses{lem:smoothOfRelativeDimension_smooth,
        lem:free_finrank_kaehler_localization,
        lem:exists_standardSmooth_chart_of_smooth,
        lem:stalk_standardSmooth_localization_of_smooth}
  % SOURCE: [Hartshorne], II.8, Theorem 8.15, p.~177 (an irreducible
  % separated finite-type scheme over an algebraically closed field is
  % nonsingular iff its sheaf of differentials is locally free of rank
  % $n = \dim X$). Read from
  % references/hartshorne-algebraic-geometry.pdf, PDF page 194 (= doc
  % p.~177), transcribed character-by-character from the rendered scanned
  % page image (the PDF has no text layer).
  % SOURCE QUOTE: "Let $X$ be an irreducible separated scheme of finite type
  % over an algebraically closed field $k$. Then $\Omega_{X/k}$ is a locally
  % free sheaf of rank $n = \dim X$ if and only if $X$ is a nonsingular
  % variety over $k$."
  \textit{Source: Hartshorne, II.8, Theorem~8.15.}
  Let \(\bar k\) be an algebraically closed field and let \(C\) be an
  integral curve of finite type over \(\bar k\) whose structure morphism
  is \(\texttt{SmoothOfRelativeDimension}\ 1\). Then the sheaf of relative
  differentials \(\Omega_{C/\bar k}\) is a locally free
  \(\mathcal O_C\)-module of rank one.
\end{lemma}

% SOURCE QUOTE PROOF: "If $x \in X$ is a closed point, then the local ring
% $B = \mathcal{O}_{X,x}$ has dimension $n$, residue field $k$, and is a
% localization of a $k$-algebra of finite type. Furthermore the module
% $\Omega_{B/k}$ of differentials of $B$ over $k$ is equal to the stalk
% $(\Omega_{X/k})_x$ of the sheaf $\Omega_{X/k}$, and we can apply (8.8) and
% we see that $(\Omega_{X/k})_x$ is free of rank $n$ if and only if $B$ is a
% regular local ring. Now the theorem follows in view of (8.14A) and
% (Ex. 5.7)."
\begin{proof}
  \uses{lem:smoothOfRelativeDimension_smooth,
        lem:exists_standardSmooth_chart_of_smooth,
        lem:stalk_standardSmooth_localization_of_smooth,
        lem:free_finrank_kaehler_localization}
  By \cref{lem:smoothOfRelativeDimension_smooth} the structure morphism is a
  smooth morphism, so \(C\) is a nonsingular variety over the algebraically
  closed field \(\bar k\) (all its local rings are regular). By Hartshorne's
  Theorem~II.8.15 the sheaf \(\Omega_{C/\bar k}\) is then locally free of
  rank \(n = \dim C = 1\).

  \emph{Stalk-level realisation.} For the codimension-one application only the
  rank-one freeness of the stalk \(\Omega_{C/\bar k, x}\) is consumed, and that
  is obtained directly from the standard-smooth chart, with no recourse to the
  global locally-free statement. By
  \cref{lem:exists_standardSmooth_chart_of_smooth} the point \(x\) lies on a
  relative-dimension-one standard-smooth affine chart
  \(\operatorname{Spec} S\) of \(C\), and by
  \cref{lem:stalk_standardSmooth_localization_of_smooth} the stalk
  \(\mathcal O_{C, x} = T = S_{\mathfrak q}\) is the localisation of \(S\) at
  the corresponding prime \(\mathfrak q\), with \(S\) standard-smooth of
  relative dimension one over \(\bar k\). Applying the pure-algebra core
  \cref{lem:free_finrank_kaehler_localization} to this localisation gives that
  \(\Omega_{T/\bar k} = \Omega_{C/\bar k, x}\) is a free
  \(\mathcal O_{C, x}\)-module of rank one with
  \(\operatorname{finrank}_{\mathcal O_{C, x}} \Omega_{C/\bar k, x} = 1\).
\end{proof}

% NOTE: the codimension-one rationality lemma
% \cref{lem:codimOne_point_residueField_eq_kbar}
% (\lean{AlgebraicGeometry.residueField_eq_of_coheight_eq_one}) and its proof now
% live in \cref{chap:RiemannRoch_ResidueFieldKbar}, the owning chapter of
% AlgebraicJacobian/RiemannRoch/ResidueFieldKbar.lean.  It is \uses{}'d by
% \cref{lem:stalk_isDVR_of_smooth} across chapters.

\begin{lemma}

\leanok
[Surjectivity of the structure map onto the residue field at a maximal localisation]
  \label{lem:algebraMap_residueField_surjective}
  \lean{Algebra.algebraMap_residueField_surjective_of_isLocalization_atPrime}
  Let \(R\) be a commutative ring and \(\mathfrak p \subseteq R\) a
  \emph{maximal} ideal. Let \(R_{\mathfrak p}\) be a localisation of \(R\) at
  \(\mathfrak p\) --- a local ring with
  \(\texttt{IsLocalization.AtPrime}\ R_{\mathfrak p}\ \mathfrak p\) --- and let
  \(\kappa(R_{\mathfrak p}) = R_{\mathfrak p} / \mathfrak m_{R_{\mathfrak p}}\)
  denote its residue field. Then the composite structure map
  \[
    R \;\longrightarrow\; R_{\mathfrak p} \;\longrightarrow\;
    \kappa(R_{\mathfrak p})
  \]
  is surjective.
\end{lemma}

\begin{proof}
  \leanok
  The residue map \(R_{\mathfrak p} \to \kappa(R_{\mathfrak p})\) is surjective, so any element of
  \(\kappa(R_{\mathfrak p})\) is the residue class of some
  \(t \in R_{\mathfrak p}\); by surjectivity of the localisation map onto
  fractions, write \(t = r / s\) with \(r \in R\) and \(s \in R \setminus
  \mathfrak p\). The composite \(R \to \kappa(R_{\mathfrak p})\) kills exactly
  \(\mathfrak p\): a class \(\overline{a}\) (for \(a \in R\)) vanishes in
  \(\kappa(R_{\mathfrak p})\) iff \(a \in \mathfrak p\).
  Because \(\mathfrak p\) is maximal and
  \(s \notin \mathfrak p\), the ideal \((s) + \mathfrak p\) strictly contains
  \(\mathfrak p\), hence equals the whole ring \(R\):
  there exist \(u \in R\) and \(q \in \mathfrak p\) with \(u s = 1 - q\), i.e.
  \(s\) is invertible modulo \(\mathfrak p\).
  Consequently \(\overline{u}\) is the inverse of \(\overline{s}\) in
  \(\kappa(R_{\mathfrak p})\), and the image of \(r u \in R\) is
  \(\overline{r}\,\overline{s}^{-1}\), which is precisely the residue class of
  \(t = r/s\), the chosen element. Thus every element of
  \(\kappa(R_{\mathfrak p})\) is hit by \(R\).
\end{proof}

\begin{lemma}

\leanok
[Cotangent space at a rational point of a smooth curve is one-dimensional]
  \label{lem:cotangentSpace_finrank_eq_one_of_smooth}
  \lean{AlgebraicGeometry.finrank_cotangentSpace_stalk_eq_one_of_smooth}
  \uses{lem:algebraMap_residueField_surjective,
        lem:finrank_cotangent_localization_eq_one,
        lem:exists_standardSmooth_chart_of_smooth,
        lem:stalk_standardSmooth_localization_of_smooth,
        lem:smooth_kaehler_locallyFree_rank_one,
        lem:finrank_cotangentSpace_eq_rank_kaehler,
        thm:krullDim_curve_le_one}
  % SOURCE: [Hartshorne], II.8, Proposition 8.7, p.~174 (for a local ring
  % $B$ containing a field $k$ isomorphic to its residue field, the conormal
  % map $\mathfrak{m}/\mathfrak{m}^2 \to \Omega_{B/k} \otimes_B k$ is an
  % isomorphism). Read from
  % references/hartshorne-algebraic-geometry.pdf, PDF page 191 (= doc
  % p.~174), transcribed character-by-character from the rendered scanned
  % page image (the PDF has no text layer).
  % SOURCE QUOTE: "Let $B$ be a local ring which contains a field $k$
  % isomorphic to its residue field $B/\mathfrak{m}$. Then the map
  % $\delta: \mathfrak{m}/\mathfrak{m}^2 \to \Omega_{B/k} \otimes_B k$ of
  % (8.4A) is an isomorphism."
  \textit{Source: Hartshorne, II.8, Proposition~8.7.}
  Let \(C\) be a smooth (\(\texttt{SmoothOfRelativeDimension}\ 1\)) integral
  curve over an algebraically closed field \(\bar k\), assumed to have Krull
  dimension at most one (\(\operatorname{krullDim} C.\mathrm{left} \le 1\)), and
  let \(x \in C\) be a codimension-one point (\(\texttt{Order.coheight}\,x = 1\)).
  Then the cotangent space of the stalk is one-dimensional over the residue
  field:
  \[
    \dim_{\kappa(x)} \mathfrak m_x / \mathfrak m_x^2
    \;=\;
    \operatorname{finrank}_{\kappa(x)}
      \bigl(\texttt{IsLocalRing.CotangentSpace}\ \mathcal O_{C, x}\bigr)
    \;=\; 1.
  \]
\end{lemma}

% SOURCE QUOTE PROOF: "According to (8.4A), the cokernel of $\delta$ is
% $\Omega_{k/k} = 0$, so $\delta$ is surjective. To show that $\delta$ is
% injective, it will be sufficient to show that the map
% $\delta': \Hom_k(\Omega_{B/k} \otimes k, k) \to
% \Hom_k(\mathfrak{m}/\mathfrak{m}^2, k)$ of dual vector spaces is
% surjective. The term on the left is isomorphic to
% $\Hom_B(\Omega_{B/k}, k)$, which by definition of the differentials, can
% be identified with the set $\Der_k(B, k)$ of $k$-derivations of $B$ to
% $k$. If $d: B \to k$ is a derivation, then $\delta'(d)$ is obtained by
% restricting to $\mathfrak{m}$, and noting that $d(\mathfrak{m}^2) = 0$.
% Now, to show that $\delta'$ is surjective, let
% $h \in \Hom(\mathfrak{m}/\mathfrak{m}^2, k)$. For any $b \in B$, we can
% write $b = \lambda + c$, $\lambda \in k$, $c \in \mathfrak{m}$, in a
% unique way. Define $db = h(\bar c)$, where $\bar c \in
% \mathfrak{m}/\mathfrak{m}^2$ is the image of $c$. Then one verifies
% immediately that $d$ is a $k$-derivation of $B$ to $k$, and that
% $\delta'(d) = h$. Thus $\delta'$ is surjective, as required."
\begin{proof}
  \uses{lem:algebraMap_residueField_surjective,
        lem:finrank_cotangent_localization_eq_one,
        lem:exists_standardSmooth_chart_of_smooth,
        lem:stalk_standardSmooth_localization_of_smooth,
        lem:finrank_cotangentSpace_eq_rank_kaehler,
        thm:krullDim_curve_le_one}
        \leanok
  Write \(B = \mathcal O_{C, x}\) for the stalk.
  By \cref{lem:exists_standardSmooth_chart_of_smooth} the point \(x\) lies on a
  relative-dimension-one standard-smooth affine chart
  \(\operatorname{Spec} S\) of \(C\), and by
  \cref{lem:stalk_standardSmooth_localization_of_smooth} the stalk
  \(B = T = S_{\mathfrak q}\) is the localisation of \(S\) at the prime
  \(\mathfrak q \subseteq S\) corresponding to \(x\), with \(S\) standard-smooth
  of relative dimension one over \(\bar k\) and the scalar tower
  \(\bar k \to S \to T\). As a standard-smooth \(\bar k\)-algebra, \(S\) is of
  finite type over \(\bar k\).

  \emph{The residue field \(\kappa(x)\) is \(\bar k\)-rational.}
  Since \(\operatorname{coheight} x = 1\) and \(\operatorname{krullDim}
  C.\mathrm{left} \le 1\), the closed subset \(\overline{\{x\}}\) cannot be
  properly specialised within \(C\), so \(x\) is a closed point of \(C\);
  the global dimension bound \(\operatorname{krullDim} C.\mathrm{left} \le 1\) is exactly
  what excludes a longer specialisation chain and is supplied by
  \cref{thm:krullDim_curve_le_one}. A closed point of the affine chart
  \(\operatorname{Spec} S\) corresponds to a \emph{maximal} ideal
  \(\mathfrak q\) of \(S\). By
  \cref{lem:algebraMap_residueField_surjective} the structure map
  \(S \to \kappa(x) = \kappa(B)\) is therefore surjective, so \(\kappa(x)\) is a
  quotient of \(S\) and hence itself of finite type over \(\bar k\). A
  field that is of finite type as an algebra over the algebraically closed
  field \(\bar k\) is equal to \(\bar k\)
  (the structure map \(\bar k \to \kappa(x)\) is bijective by algebraic closedness).
  In particular
  \(\kappa(x)\) is \(\bar k\)-rational, which supplies the two residue-field
  vanishings \(\Omega_{\kappa(x)/\bar k} = 0\) and
  \(H^1_{\mathrm{cot}}(\bar k, \kappa(x)) = 0\) consumed below.

  Applying these data, together with
  the residue-field vanishings above, to the pure-algebra core
  \cref{lem:finrank_cotangent_localization_eq_one} yields directly
  \[
    \operatorname{finrank}_{\kappa(x)}
      \bigl(\texttt{IsLocalRing.CotangentSpace}\ B\bigr) = 1,
    \qquad\text{i.e.}\qquad
    \dim_{\kappa(x)} \mathfrak m_x / \mathfrak m_x^2 = 1.
  \]
  \cref{lem:finrank_cotangent_localization_eq_one} is the conormal
  comparison \cref{lem:finrank_cotangentSpace_eq_rank_kaehler} of
  \cref{chap:RiemannRoch_SmoothRegular} (the numerical corollary of the
  isomorphism
  \(\mathfrak m_x/\mathfrak m_x^2 \cong \Omega_{B/\bar k} \otimes_B \kappa(x)\),
  Hartshorne~II.8.7) applied to the rank-one freeness of
  \(\Omega_{B/\bar k}\).
\end{proof}

\begin{lemma}

\leanok
[Smooth of relative dimension one \(\Rightarrow\) DVR stalk at a
codimension-one point]
  \label{lem:stalk_isDVR_of_smooth}
  \lean{AlgebraicGeometry.isDiscreteValuationRing_stalk_of_smooth}
  \uses{lem:cotangentSpace_finrank_eq_one_of_smooth,
        lem:isDVR_stalk_of_finrank_cotangentSpace,
        lem:isDVR_localization_standardSmooth,
        lem:exists_standardSmooth_chart_of_smooth,
        lem:stalk_standardSmooth_localization_of_smooth}
  Let \(\bar k\) be an algebraically closed field and let \(C\) be an
  integral curve over \(\bar k\) whose structure morphism is
  \(\texttt{SmoothOfRelativeDimension}\ 1\). Then for every codimension-one
  point \(x \in C\) (\(\texttt{Order.coheight}\,x = 1\)) the stalk
  \(\mathcal O_{C, x}\) is a discrete valuation ring
  (\(\texttt{IsDiscreteValuationRing}\,\mathcal O_{C, x}\)).
\end{lemma}

\begin{proof}
  \uses{lem:exists_standardSmooth_chart_of_smooth,
        lem:stalk_standardSmooth_localization_of_smooth,
        lem:isDVR_localization_standardSmooth,
        lem:codimOne_point_residueField_eq_kbar}
  Let \(x \in C\) be a codimension-one point, and write
  \(T = \mathcal O_{C, x}\) for the stalk. As the local ring of a point on an
  integral, locally Noetherian curve, \(T\) is a Noetherian local integral
  domain. Since \(x\) is a codimension-one point of an integral curve over the
  algebraically closed field \(\bar k\), it is closed with residue field
  \(\kappa(x) = \bar k\) (\cref{lem:codimOne_point_residueField_eq_kbar}), so the
  residue field is \(\bar k\)-rational.

  By \cref{lem:exists_standardSmooth_chart_of_smooth} the point \(x\) lies on a
  relative-dimension-one standard-smooth affine chart
  \(\operatorname{Spec} S\) of \(C\), and by
  \cref{lem:stalk_standardSmooth_localization_of_smooth} the stalk
  \(T = S_{\mathfrak q}\) is the localisation of \(S\) at the corresponding
  prime, with \(S\) standard-smooth of relative dimension one over \(\bar k\)
  and the localisation and base-field structures compatible. These are exactly
  the hypotheses of the pure-algebra core
  \cref{lem:isDVR_localization_standardSmooth}: a localisation \(T\) of a
  relative-dimension-one standard-smooth \(\bar k\)-algebra that is a Noetherian
  local integral domain with \(\bar k\)-rational residue field is a discrete
  valuation ring. Applying it to \(T\) yields
  \(\texttt{IsDiscreteValuationRing}\,\mathcal O_{C, x}\).
\end{proof}

\begin{lemma}

\leanok
[Codimension one \(\Rightarrow\) Krull dimension one of the stalk]
  \label{lem:coheight_eq_ringKrullDim_stalk}
  \lean{AlgebraicGeometry.ringKrullDim_stalk_eq_one_of_coheight_eq_one}
  % SOURCE: [Stacks Project], Tag 02IZ = Lemma 28.10.3
  % (`properties-lemma-codimension-local-ring`, Chapter 28 ``Properties of
  % Schemes''): the codimension of an irreducible closed subset equals the
  % Krull dimension of the local ring at its generic point. Read from
  % references/stacks-coheight-properties.tex, lines 1258--1274.
  % SOURCE QUOTE: "Let $X$ be a scheme. Let $Y \subset X$ be an irreducible
  % closed subset. Let $\xi \in Y$ be the generic point. Then
  % $$
  % \text{codim}(Y, X) = \dim(\mathcal{O}_{X, \xi})
  % $$
  % where the codimension is as defined in
  % Topology, Definition \ref{topology-definition-codimension}."
  \textit{Source: Stacks Project, Tag 02IZ (Lemma 28.10.3): for an
  irreducible closed subset \(Y\) with generic point \(\xi\),
  \(\operatorname{codim}(Y, X) = \dim \mathcal O_{X, \xi}\).}
  Let \(C\) be an integral scheme and let \(x \in C\) be a point with
  \(\texttt{Order.coheight}\,x = 1\), i.e.\ the irreducible closed subset
  \(\overline{\{x\}}\) (of which \(x\) is the generic point) has
  codimension one in \(C\). Then \(\operatorname{ringKrullDim}
  \mathcal O_{C, x} = 1\).
\end{lemma}

\begin{proof}
  The point \(x\) is the generic point of the irreducible closed subset
  \(Y := \overline{\{x\}}\). By the Stacks Project (Tag 02IZ) the
  codimension of \(Y\) in \(C\) equals the Krull dimension of the local ring
  at its generic point:
  \(\operatorname{codim}(Y, C) = \dim \mathcal O_{C, x}\). The topological
  coheight \(\texttt{Order.coheight}\,x\) in the specialisation order on
  \(C\) is, by definition, the codimension \(\operatorname{codim}(Y, C)\).
  The hypothesis \(\texttt{Order.coheight}\,x = 1\) therefore yields
  \(\operatorname{ringKrullDim} \mathcal O_{C, x} = 1\).
\end{proof}

\begin{lemma}[Cotangent-space characterisation of DVRs]
  \label{lem:finrank_cotangentSpace_eq_one_iff_isDVR}
  \lean{IsLocalRing.finrank_CotangentSpace_eq_one_iff}
  \mathlibok
  \textit{Provided by Mathlib
  (\texttt{Mathlib.RingTheory.DiscreteValuationRing.TFAE}).}
  For a Noetherian local integral domain \(R\), the residue-field dimension
  of the cotangent space satisfies
  \(\dim_{\kappa(R)} \mathfrak m_R / \mathfrak m_R^2 = 1\) if and only if
  \(R\) is a discrete valuation ring. (This is one clause of the standard
  \texttt{TFAE} for discrete valuation rings,
  \(\texttt{IsDiscreteValuationRing.TFAE}\), exposed as the stand-alone
  equivalence \(\texttt{finrank\_CotangentSpace\_eq\_one\_iff}\).)
\end{lemma}

\begin{lemma}

\leanok
[Scheme-stalk wrapper: one-dimensional cotangent space \(\Rightarrow\) DVR
stalk]
  \label{lem:isDVR_stalk_of_finrank_cotangentSpace}
  \lean{AlgebraicGeometry.isDiscreteValuationRing_stalk_of_finrank_cotangentSpace_eq_one}
  \uses{lem:finrank_cotangentSpace_eq_one_iff_isDVR}
  Let \(X\) be an integral, locally Noetherian scheme and let \(x \in X\).
  If the cotangent space of the stalk is one-dimensional over the residue
  field,
  \[
    \operatorname{finrank}_{\kappa(x)}
      \bigl(\texttt{IsLocalRing.CotangentSpace}\ \mathcal O_{X, x}\bigr)
    \;=\; 1,
  \]
  then the stalk \(\mathcal O_{X, x}\) is a discrete valuation ring
  (\(\texttt{IsDiscreteValuationRing}\,\mathcal O_{X, x}\)).
\end{lemma}

\begin{proof}
  \uses{lem:finrank_cotangentSpace_eq_one_iff_isDVR}
  The stalk \(\mathcal O_{X, x}\) of an integral, locally Noetherian scheme
  at a point is a Noetherian local integral domain: it is a local ring
  (\(\texttt{IsLocalRing}\)), Noetherian (\(\texttt{IsNoetherianRing}\), from
  local Noetherianity of \(X\)), and an integral domain
  (\(\texttt{IsDomain}\), from integrality of \(X\)). These three properties
  are precisely the hypotheses of \cref{lem:finrank_cotangentSpace_eq_one_iff_isDVR}
  (Mathlib's \(\texttt{IsLocalRing.finrank\_CotangentSpace\_eq\_one\_iff}\)),
  whose forward direction converts the hypothesis
  \(\operatorname{finrank}_{\kappa(x)}
  (\texttt{CotangentSpace}\ \mathcal O_{X, x}) = 1\) into
  \(\texttt{IsDiscreteValuationRing}\,\mathcal O_{X, x}\). This lemma reduces
  the discrete valuation ring conclusion to the single numerical condition
  \(\operatorname{finrank}_{\kappa(x)} \mathfrak m_x/\mathfrak m_x^2 = 1\).
\end{proof}

\begin{lemma}

\leanok
[Regular local ring of Krull dimension one \(\Rightarrow\) DVR]
  \label{lem:regularLocal_dim_one_isDVR}
  \lean{AlgebraicGeometry.isDiscreteValuationRing_of_isRegularLocalRing_of_ringKrullDim_eq_one}
  \uses{lem:finrank_cotangentSpace_eq_one_iff_isDVR}
  % SOURCE: [Stacks Project], Tag 00PD = Lemma 10.119.7
  % (`algebra-lemma-characterize-dvr`, Section 10.119 ``Around
  % Krull--Akizuki''), the equivalence of ``regular local ring of
  % dimension one'' and ``discrete valuation ring''. Read from
  % references/stacks-regular-dvr-algebra.tex, lines 29107--29121.
  % SOURCE QUOTE: "Let $A$ be a ring. The following are equivalent.
  % \begin{enumerate}
  % \item The ring $A$ is a discrete valuation ring.
  % \item The ring $A$ is a valuation ring and Noetherian but not a field.
  % \item The ring $A$ is a regular local ring of dimension $1$.
  % \item The ring $A$ is a Noetherian local domain with maximal ideal
  % $\mathfrak m$ generated by a single nonzero element.
  % \item The ring $A$ is a Noetherian local normal domain of dimension $1$.
  % \end{enumerate}"
  \textit{Source: Stacks Project, Tag 00PD (Lemma 10.119.7), the
  equivalence (1)\,\(\Leftrightarrow\)\,(3): a regular local ring of
  dimension one is a discrete valuation ring.}
  Let \(R\) be a Noetherian local integral domain that is a regular local
  ring (\(\texttt{IsRegularLocalRing}\,R\)) with
  \(\operatorname{ringKrullDim} R = 1\). Then \(R\) is a discrete valuation
  ring (\(\texttt{IsDiscreteValuationRing}\,R\)).
\end{lemma}

\begin{proof}
  \uses{lem:finrank_cotangentSpace_eq_one_iff_isDVR}
  This is the implication (3)\,\(\Rightarrow\)\,(1) of the Stacks
  characterisation quoted above, via the cotangent-space route: by the
  definition of a regular local ring, the residue-field dimension of the
  cotangent space \(\mathfrak m_R / \mathfrak m_R^2\) equals the Krull
  dimension of \(R\), which is \(1\) by hypothesis; hence
  \(\dim_{\kappa(R)} \mathfrak m_R / \mathfrak m_R^2 = 1\), and
  \cref{lem:finrank_cotangentSpace_eq_one_iff_isDVR}
  converts this into \(\texttt{IsDiscreteValuationRing}\,R\).
\end{proof}

\begin{theorem}

\leanok
[Smooth of relative dimension one \(\Rightarrow\) regular in codimension one]
  \label{thm:smooth_isRegularInCodimOne}
  \lean{AlgebraicGeometry.instIsRegularInCodimensionOneOfSmooth}
  \uses{def:isRegularInCodimensionOne, lem:stalk_isDVR_of_smooth,
        thm:krullDim_curve_le_one}
  Let \(\bar k\) be an algebraically closed field and let \(C\) be an
  integral curve over
  \(\bar k\) whose structure morphism
  \(C.\mathrm{hom} : C \to \operatorname{Spec} \bar k\) is
  \(\texttt{SmoothOfRelativeDimension}\ 1\). Then the underlying scheme
  \(C.\mathrm{left}\) is \emph{regular in codimension one}
  (\cref{def:isRegularInCodimensionOne}): for every point \(x \in C\) of
  codimension one (\(\texttt{Order.coheight}\,x = 1\)) the local ring
  \(\mathcal O_{C, x}\) is a discrete valuation ring.
\end{theorem}

\begin{proof}
  \uses{def:isRegularInCodimensionOne, lem:stalk_isDVR_of_smooth,
        thm:krullDim_curve_le_one}
  Let \(x \in C\) be a point with \(\texttt{Order.coheight}\,x = 1\). By
  \cref{lem:stalk_isDVR_of_smooth}, applied to the
  smooth-of-relative-dimension-one structure morphism of \(C\), the stalk
  \(\mathcal O_{C, x}\) is a discrete valuation ring, which is exactly the
  per-point datum required by \cref{def:isRegularInCodimensionOne}. The
  per-point DVR discharge routes through the one-dimensional cotangent space
  \cref{lem:cotangentSpace_finrank_eq_one_of_smooth}, whose
  \(\operatorname{krullDim} C.\mathrm{left} \le 1\) hypothesis is the global
  curve-dimension bound \cref{thm:krullDim_curve_le_one}.

  As \(x\) was an arbitrary codimension-one point, \(C.\mathrm{left}\) is
  regular in codimension one. This is the abstract-curve generalisation of
  \cref{thm:projectiveLineBar_isRegularInCodimOne}, which establishes the same
  conclusion via the explicit two-chart affine cover of \(\mathbb P^1_{\bar k}\)
  (each chart an affine line, hence a principal ideal domain, whose localisation
  at a height-one prime is a discrete valuation ring).
\end{proof}

\section{Carrier substrate for \(\mathcal O_C(D)\)}

This section assembles the substrate helpers for \cref{def:sheafOf}.
They generalise the closed-point carrier machinery of
\cref{chap:RiemannRoch_OCofP} from a single closed point \(P\) to an
arbitrary Weil divisor \(D\). Throughout, \(C\) is integral, locally
Noetherian and regular in codimension one over \(\bar k\), so that the
local ring \(\mathcal O_{C, Q}\) at each prime divisor \(Q\) is a
discrete valuation ring and \(\ord_Q\) is its valuation.

\begin{definition}

\leanok
[Carrier set of \(\mathcal O_C(D)\) over an open]
  \label{def:sheafOf_carrierSet}
  \lean{AlgebraicGeometry.Scheme.WeilDivisor.sheafOf.carrierSet}
  \uses{def:codim1_cycles, def:prime_divisor, def:order_at_point}
  For a Weil divisor \(D : C.\text{left.PrimeDivisor} \to_0 \mathbb Z\) and
  an open \(U \subseteq C\), the \emph{carrier set}
  \(\mathrm{carrierSet}_D(U) \subseteq K(C)\) is
  \[
    \mathrm{carrierSet}_D(U) \;=\;
    \bigl\{\, f \in K(C) \;\big|\;
      f = 0 \,\lor\,
      \bigl(\forall Q,\ Q.\mathrm{point} \in U \;\Rightarrow\;
        -\,D(Q) \le \ord_Q(f)\bigr) \,\bigr\}.
  \]
  The disjunctive shape \(f = 0 \lor \dots\) (in contrast with the pure
  conjunction used in \texttt{RR.3}'s closed-point carrier
  \cref{def:lineBundleAtClosedPoint_carrierSet}) is \emph{required}: for a
  prime divisor \(Q\) at which \(D(Q) < 0\) the bound \(-D(Q)\) is strictly
  positive, so the junk value \(\ord_Q(0) = 0\) need not satisfy
  \(-D(Q) \le \ord_Q(0)\), and the zero function must be admitted by a
  separate clause for the carrier to be a submodule. It is the underlying
  set of the section module of \(\mathcal O_C(D)\) over \(U\).
\end{definition}

\begin{lemma}

\leanok
[Monotonicity of the carrier set]
  \label{lem:sheafOf_carrierSet_mono}
  \lean{AlgebraicGeometry.Scheme.WeilDivisor.sheafOf.carrierSet_mono}
  \uses{def:sheafOf_carrierSet}
  If \(V \subseteq U\) are opens of \(C\) then
  \(\mathrm{carrierSet}_D(U) \subseteq \mathrm{carrierSet}_D(V)\): the order
  conditions over the smaller open \(V\) range over a subset of the prime
  divisors involved over \(U\) and are therefore implied by them, while the
  \(f = 0\) clause is preserved verbatim. This antitone inclusion supplies
  the restriction map of the carrier presheaf.
\end{lemma}

\begin{proof}
  \uses{def:sheafOf_carrierSet}
  A function \(f \in \mathrm{carrierSet}_D(U)\) either is \(0\), in which
  case it lies in \(\mathrm{carrierSet}_D(V)\) by the same \(f = 0\) clause,
  or satisfies \(-D(Q) \le \ord_Q(f)\) for every \(Q\) with
  \(Q.\mathrm{point} \in U\); since \(V \subseteq U\), the hypothesis
  \(Q.\mathrm{point} \in V\) entails \(Q.\mathrm{point} \in U\), so the same
  bound holds for every \(Q\) with \(Q.\mathrm{point} \in V\). Hence
  \(f \in \mathrm{carrierSet}_D(V)\).
\end{proof}

\begin{definition}

\leanok
[Carrier submodule of \(\mathcal O_C(D)\) over an open]
  \label{def:sheafOf_carrierSubmodule}
  \lean{AlgebraicGeometry.Scheme.WeilDivisor.sheafOf.carrierSubmodule}
  \uses{def:sheafOf_carrierSet, def:order_at_point}
  The carrier set \(\mathrm{carrierSet}_D(U)\) is a \(\bar k\)-submodule of
  \(K(C)\). It contains \(0\) (the \(f = 0\) clause). It is closed under
  addition: if \(a, b\) are both nonzero with \(a + b \neq 0\), then for
  every prime divisor \(Q\) the non-archimedean inequality
  \(\min(\ord_Q a, \ord_Q b) \le \ord_Q(a + b)\) combines with the two hypotheses
  \(-D(Q) \le \ord_Q a\) and \(-D(Q) \le \ord_Q b\) to give
  \(-D(Q) \le \ord_Q(a + b)\); the degenerate cases where one summand or the
  sum is \(0\) fall back to the \(f = 0\) clause. It is closed under scaling
  by a constant \(c \in \bar k\): for \(c \neq 0\) the image of \(c\) in
  each stalk \(\mathcal O_{C, Q}\) is a unit, hence has order \(\ge 0\), so
  \(\ord_Q(c \cdot f) = \ord_Q(c) + \ord_Q(f) \ge \ord_Q(f) \ge -D(Q)\),
  while \(c = 0\) gives \(c \cdot f = 0\).
\end{definition}

\begin{definition}

\leanok
[Bot-trivialisation submodule]
  \label{def:sheafOf_trivAtBot}
  \lean{AlgebraicGeometry.Scheme.WeilDivisor.sheafOf.trivAtBot}
  For an open \(U \subseteq C\), the \(\bar k\)-submodule
  \(\mathrm{trivAtBot}(U) \subseteq K(C)\) equals all of \(K(C)\) when
  \(U \ne \emptyset\) and equals \(\{0\}\) when \(U = \emptyset\); concretely
  \(\mathrm{trivAtBot}(U) = \{\, f \mid U \ne \emptyset \,\lor\, f = 0\,\}\).
  Intersecting the carrier submodule with this factor enforces the
  sheaf-at-\(\emptyset\) condition \(\Gamma(\emptyset, \mathcal O_C(D)) = 0\)
  required by the Grothendieck topology of opens (whose empty sieve covers
  \(\emptyset\), forcing the section module there to be terminal).
\end{definition}

\begin{definition}

\leanok
[Sheaf-corrected carrier submodule enforcing \(F(\bot) = 0\)]
  \label{def:sheafOf_carrierSubmoduleSheaf}
  \lean{AlgebraicGeometry.Scheme.WeilDivisor.sheafOf.carrierSubmoduleSheaf}
  \uses{def:sheafOf_carrierSubmodule, def:sheafOf_trivAtBot}
  The \emph{sheaf-corrected carrier submodule} is the meet
  \[
    \mathrm{carrierSubmoduleSheaf}_D(U) \;:=\;
    \mathrm{carrierSubmodule}_D(U) \;\sqcap\; \mathrm{trivAtBot}(U).
  \]
  On a nonempty open it equals \(\mathrm{carrierSubmodule}_D(U)\) (then
  \(\mathrm{trivAtBot}(U) = K(C)\)); at \(U = \emptyset\) it equals \(\{0\}\)
  (the meet annihilates, since \(\mathrm{trivAtBot}(\emptyset) = \{0\}\)).
  This is the corrected substrate consumed by the carrier presheaf, exactly
  as in \texttt{RR.3}'s
  \cref{def:lineBundleAtClosedPoint_carrierSubmoduleSheaf}.
\end{definition}

\begin{lemma}

\leanok
[Monotonicity of the sheaf-corrected carrier submodule]
  \label{lem:sheafOf_carrierSubmoduleSheaf_le}
  \lean{AlgebraicGeometry.Scheme.WeilDivisor.sheafOf.carrierSubmoduleSheaf_le}
  \uses{def:sheafOf_carrierSubmoduleSheaf, lem:sheafOf_carrierSet_mono}
  If \(V \subseteq U\) are opens with \(V \ne \emptyset\), then
  \(\mathrm{carrierSubmoduleSheaf}_D(U) \le
  \mathrm{carrierSubmoduleSheaf}_D(V)\). The hypothesis \(V \ne \emptyset\)
  is necessary: a nonzero section over \(U\) need not satisfy the
  trivialisation-at-\(\emptyset\) condition. This inclusion supplies the
  restriction maps of the carrier presheaf on nonempty targets.
\end{lemma}

\begin{proof}
  \uses{lem:sheafOf_carrierSet_mono, def:sheafOf_trivAtBot}
  Let \(x \in \mathrm{carrierSubmoduleSheaf}_D(U)\), so
  \(x \in \mathrm{carrierSet}_D(U)\) and \(x \in \mathrm{trivAtBot}(U)\). By
  \cref{lem:sheafOf_carrierSet_mono} and \(V \subseteq U\),
  \(x \in \mathrm{carrierSet}_D(V)\); and since \(V \ne \emptyset\),
  \(x \in \mathrm{trivAtBot}(V)\) holds by the left disjunct. Hence
  \(x \in \mathrm{carrierSubmoduleSheaf}_D(V)\).
\end{proof}

\begin{definition}

\leanok
[Type-level subfunctor presentation of the carrier]
  \label{def:sheafOf_carrierTypeSubfunctor}
  \lean{AlgebraicGeometry.Scheme.WeilDivisor.sheafOf.carrierTypeSubfunctor}
  \uses{def:sheafOf_carrierSubmoduleSheaf}
  The carrier of \(\mathcal O_C(D)\) is presented as a subfunctor of the
  type-valued presheaf \(U \mapsto (U \to K(C))\) of arbitrary
  \(K(C)\)-valued functions: a section over \(U\) is a constant function
  with value \(f \in \mathrm{carrierSubmoduleSheaf}_D(U)\). On a restriction
  \(V \subseteq U\), a section with value \(f\) maps to the section with
  value \(f\) when \(V \ne \emptyset\) (membership transported by
  \cref{lem:sheafOf_carrierSubmoduleSheaf_le}) and to the unique section
  over \(V = \emptyset\) otherwise. This presentation reduces the sheaf
  condition to a stalk-locality check against the ambient sheaf of
  type-valued functions, exploiting that any two nonempty opens of the
  irreducible scheme \(C\) intersect.
\end{definition}

\begin{definition}

\leanok
[The carrier presheaf of \(\mathcal O_C(D)\)]
  \label{def:sheafOf_carrierPresheaf}
  \lean{AlgebraicGeometry.Scheme.WeilDivisor.sheafOf.carrierPresheaf}
  \uses{def:sheafOf_carrierSubmoduleSheaf, lem:sheafOf_carrierSubmoduleSheaf_le}
  The sheaf-corrected carrier submodules assemble into a
  \(\bar k\)-linear functor
  \[
    \texttt{carrierPresheaf}_D : (\mathrm{Opens}\,|C|)^{\mathrm{op}}
    \longrightarrow \mathbf{Mod}_{\bar k},
    \qquad U \longmapsto \mathrm{carrierSubmoduleSheaf}_D(U),
  \]
  whose restriction along \(V \subseteq U\) is the zero map when
  \(V = \emptyset\) and the inclusion
  \(\mathrm{carrierSubmoduleSheaf}_D(U) \hookrightarrow
  \mathrm{carrierSubmoduleSheaf}_D(V)\) of
  \cref{lem:sheafOf_carrierSubmoduleSheaf_le} otherwise. The functor laws
  (identity and composition) are checked by cases on whether the target
  open is empty: on nonempty targets every restriction is an inclusion that
  is the identity on \(K(C)\) at the carrier level, and at the empty target
  the section module is \(\{0\}\), so both composites agree.
\end{definition}

\begin{lemma}[Carrier-presheaf restrictions are the identity on rational functions]
  \label{lem:carrierPresheaf_map_coe}
  \lean{AlgebraicGeometry.Scheme.WeilDivisor.sheafOf.carrierPresheaf_map_coe}
  \uses{def:sheafOf_carrierPresheaf}
  For every inclusion of opens \(V \subseteq U\) with non-empty target, the
  restriction map of \(\texttt{carrierPresheaf}_D\) (\cref{def:sheafOf_carrierPresheaf})
  acts as the identity on the underlying rational function: if \(x \in
  \mathrm{carrierSubmoduleSheaf}_D(U)\) restricts to \(x' \in
  \mathrm{carrierSubmoduleSheaf}_D(V)\), then \(x\) and \(x'\) are the same
  element of \(K(C)\).
\end{lemma}

\begin{proof}
  \uses{def:sheafOf_carrierPresheaf}
  By \cref{def:sheafOf_carrierPresheaf}, at a non-empty target the restriction
  map is the submodule inclusion
  \(\mathrm{carrierSubmoduleSheaf}_D(U) \hookrightarrow
  \mathrm{carrierSubmoduleSheaf}_D(V)\) (the empty-target branch is excluded by
  the hypothesis). A submodule inclusion is, by construction, the identity on
  underlying elements, so the underlying \(K(C)\)-element is preserved.
\end{proof}

\begin{lemma}

\leanok
[Sheaf property of the carrier presheaf]
  \label{lem:sheafOf_carrierPresheaf_isSheaf}
  \lean{AlgebraicGeometry.Scheme.WeilDivisor.sheafOf.carrierPresheaf_isSheaf}
  \uses{def:sheafOf_carrierPresheaf, def:sheafOf_carrierSubmoduleSheaf,
        def:sheafOf_carrierTypeSubfunctor, def:sheafOf_trivAtBot,
        def:prime_divisor, def:order_at_point}
  The presheaf \(\texttt{carrierPresheaf}_D\) is a sheaf for the
  Grothendieck topology of opens on \(C\). The proof verifies the
  unique-gluing form of the sheaf condition, discharged in two
  regimes:
  \begin{enumerate}
    \item \emph{Empty-cover regime.} When the cover sups to \(\emptyset\)
      (in particular for the empty sieve covering \(\bot\)), every member
      open is \(\emptyset\), each local section module is \(\{0\}\) by the
      \(\mathrm{trivAtBot}\) factor of
      \cref{def:sheafOf_carrierSubmoduleSheaf}, and the unique gluing is
      \(0\); existence and uniqueness are then immediate.
    \item \emph{Irreducible-gluing regime.} When the cover has a nonempty
      member, pick one, \(U_{i_0}\), and let \(v \in K(C)\) be the value of
      its section. Any two nonempty opens of the irreducible space
      underlying the integral scheme \(C\) intersect in a nonempty open, so
      on each pairwise overlap the compatibility condition — which, on
      nonempty opens, is equality of the constant \(K(C)\)-values — forces
      every nonempty member's section to carry the \emph{same} value \(v\).
      The candidate glued section is the constant function \(v\). Its
      membership in \(\mathrm{carrierSubmoduleSheaf}_D\) over the union is
      checked prime divisor by prime divisor: a prime divisor \(Q\) in the
      union lies in some member \(U_i\); if \(Q.\mathrm{point} \in U_i\)
      with \(U_i \ne \emptyset\) then \(v\) inherits the order bound
      \(-D(Q) \le \ord_Q(v)\) from the section over \(U_i\) (and the
      \(f = 0\) clause is transported when \(v = 0\)). The order conditions
      being stalk-local at each prime divisor, gluing preserves them;
      uniqueness follows because the restriction maps are the identity on
      \(K(C)\) on nonempty opens.
  \end{enumerate}
  The argument is the divisor-indexed generalisation of
  \cref{lem:lineBundleAtClosedPoint_carrierPresheaf_isSheaf}.
\end{lemma}

\section{Immediate corollaries}

\begin{lemma}\leanok
[Sheaf of the zero divisor is the structure sheaf]
  \label{lem:sheafOf_zero}
  \lean{AlgebraicGeometry.Scheme.WeilDivisor.sheafOf_zero}
  \uses{def:sheafOf}

  Let \(C\) be a smooth proper geometrically irreducible curve over
  \(\bar k\). Then the invertible sheaf associated to the zero divisor
  \(0 \in \mathrm{Div}(C)\) is canonically the structure sheaf:
  \[
    \mathcal O_C(0) \;=\; \mathcal O_C
  \]
  as sub-\(\mathcal O_C\)-modules of \(\mathcal K_C\), equivalently as
  objects in \(\mathbf{Sh}(C, \mathbf{Mod}_{\bar k})\).
\end{lemma}

\begin{proof}
  
  At \(D = 0\) the coefficient \(n_Q = 0\) at every prime divisor \(Q\), so the
  section condition of \cref{def:sheafOf} reduces, on every open
  \(U \subseteq C\), to
  \[
    \Gamma\bigl(U, \mathcal O_C(0)\bigr)
    \;=\; \{ 0 \} \cup \{ f \in K(C)^{\times} \;|\;
      \ord_Q(f) \geq 0 \text{ for every prime divisor } Q \in U \}.
  \]
  By the local description of the structure sheaf on the integral
  scheme \(C\) (the regular functions on \(U\) are exactly the rational
  functions whose order is non-negative at every prime divisor of
  \(C\) landing in \(U\); this is the standard ``non-negative order \(=\)
  regular'' identification, e.g.\ Hartshorne~II~\S6 immediately before
  Proposition~6.11), the right-hand side is precisely
  \(\Gamma(U, \mathcal O_C)\). The restriction maps on both sides are
  the identity on \(K(C)\), so the equality of presheaves promotes to an
  equality of sub-\(\mathcal O_C\)-modules of \(\mathcal K_C\).
\end{proof}

\begin{lemma}

\leanok
[Sheaf at a single closed point is isomorphic to the line bundle of
\cref{def:lineBundleAtClosedPoint}]
  \label{lem:sheafOf_singlePoint}
  \lean{AlgebraicGeometry.Scheme.WeilDivisor.sheafOf_singlePoint_iso}
  \uses{def:sheafOf, def:divisor_closed_point, def:prime_divisor,
        def:lineBundleAtClosedPoint,
        lem:lineBundleAtClosedPoint_globalSections_iff}

  Let \(C\) be a smooth proper geometrically irreducible curve over
  \(\bar k\) and let \(P \in C\) be a closed point, viewed as a Weil divisor
  \([P] \in \mathrm{Div}(C)\) via \cref{def:divisor_closed_point}. Then the
  invertible sheaf \(\mathcal O_C([P])\) of \cref{def:sheafOf} is
  \emph{isomorphic} to the line bundle of the closed point
  \(\mathcal O_C(P)\) of \cref{def:lineBundleAtClosedPoint}: there is an
  isomorphism of sheaves of \(\mathcal O_C\)-modules
  \[
    \mathcal O_C(D)\bigl|_{D = [P]}
    \;\xrightarrow{\ \sim\ }\;
    \texttt{lineBundleAtClosedPoint}_{P}
  \]
  that is the identity on the common carrier \(K(C)\); equivalently, the
  type of such isomorphisms is inhabited,
  \(\mathrm{Nonempty}\bigl(\mathcal O_C([P]) \cong
  \texttt{lineBundleAtClosedPoint}_P\bigr)\) in
  \(\mathbf{Sh}(C, \mathbf{Mod}_{\bar k})\).

  A strict sheaf \emph{equality}
  \(\mathcal O_C(D)|_{D=[P]} = \texttt{lineBundleAtClosedPoint}_P\) is not
  available: the right-hand side is assembled from the carrier data of
  \cref{chap:RiemannRoch_OCofP}, which is encapsulated and cannot be named
  for a definitional equality. The comparison is therefore stated and proved
  at the level of an isomorphism, which is all the downstream
  \(\chi\)-identity and degree arguments require.
\end{lemma}

\begin{proof}
  \uses{lem:lineBundleAtClosedPoint_globalSections_iff}
  Unfold both definitions. On the \texttt{sheafOf} side, the formal sum
  \([P]\) has coefficients \(n_Q = 1\) if \(Q = P\) and \(n_Q = 0\) otherwise
  (where \(P\) is identified with the prime divisor it determines by
  \cref{def:divisor_closed_point}). Thus
  \(\Gamma(U, \mathcal O_C([P]))\) contains \(0\) and those
  \(f \in K(C)^{\times}\) satisfying \(\ord_Q(f) \geq 0\) for every prime
  divisor \(Q \in U\) with \(Q \neq P\), plus \(\ord_P(f) \geq -1\) if
  \(P \in U\).

  On the \cref{def:lineBundleAtClosedPoint} side, the same conditions
  are exactly the ones extracted by
  \cref{lem:lineBundleAtClosedPoint_globalSections_iff} of
  \texttt{RR.3}, at every open \(U\) (not only at the global section
  \(U = C\)): a rational function \(f \in K(C)\) is a section of
  \(\mathcal O_C(P)\) on \(U\) iff it lies in the structure sheaf at every
  \(Q \in U \setminus \{P\}\) and has at most a simple pole at \(P\) when
  \(P \in U\).

  Hence over every open \(U\) the two section \(\bar k\)-submodules of
  \(K(C)\) coincide as subsets,
  \(\Gamma(U, \mathcal O_C([P])) = \Gamma(U, \texttt{lineBundleAtClosedPoint}_P)\).
  Since restriction maps on both sides are the identity on \(K(C)\) and the
  \(\bar k\)-module structure on each side is the restriction of that on
  \(K(C)\), the section-wise identity maps assemble into a morphism of
  sheaves of \(\mathcal O_C\)-modules that is the identity on the carrier;
  its section-wise inverse (again the identity on \(K(C)\)) is likewise a
  sheaf morphism. The two are mutually inverse, so they constitute an
  isomorphism \(\mathcal O_C([P]) \cong \texttt{lineBundleAtClosedPoint}_P\),
  identity on \(K(C)\). This exhibits the required inhabitant of
  \(\mathrm{Nonempty}\bigl(\mathcal O_C([P]) \cong
  \texttt{lineBundleAtClosedPoint}_P\bigr)\).
\end{proof}

\begin{lemma}
[Ideal-sheaf short exact sequence at a closed point]
  \label{lem:idealSheaf_ses_closedPoint}
  \lean{AlgebraicGeometry.Scheme.WeilDivisor.idealSheaf_ses_closedPoint}
  \uses{lem:sheafOf_ses_single_add, def:sheafOf, def:divisor_closed_point,
        def:prime_divisor}
  Let \(C\) be a smooth proper geometrically irreducible curve over
  \(\bar k\) and let \(P \in C\) be a closed point. Viewing \(P\) as a
  locally principal closed subscheme of \(C\), its structure sheaf is the
  skyscraper sheaf \(k(P)\) at \(P\) with stalk \(\bar k\), and its ideal
  sheaf is \(\mathcal O_C(-[P])\) (the sub-\(\mathcal O_C\)-module of
  \(\mathcal K_C\) locally generated by a uniformiser \(\pi_P\) at \(P\) and
  by \(1\) away from \(P\)). These fit into a short exact sequence
  \[
    0 \;\longrightarrow\; \mathcal O_C(-[P])
    \;\longrightarrow\; \mathcal O_C
    \;\longrightarrow\; k(P)
    \;\longrightarrow\; 0
  \]
  in \(\mathbf{Sh}(C, \mathbf{Mod}_{\bar k})\).
\end{lemma}

\begin{proof}
  \uses{lem:sheafOf_ses_single_add, lem:sheafOf_zero,
        lem:cokernel_sheafOf_single_add_iso_skyscraper}
  This is the inductive-step sequence \cref{lem:sheafOf_ses_single_add}
  specialised at \(D = -[P]\). Then \(\mathcal O_C(D) = \mathcal O_C(-[P])\),
  and \(D + [P] = 0\), so \(\mathcal O_C(D+[P]) = \mathcal O_C(0) =
  \mathcal O_C\) by \cref{lem:sheafOf_zero}. The inclusion
  \(\mathcal O_C(-[P]) \hookrightarrow \mathcal O_C\) is the inclusion of
  subsheaves of \(\mathcal K_C\): a section of \(\mathcal O_C(-[P])\) over
  \(U\) is a rational function vanishing to order \(\geq 1\) at \(P\) (when
  \(P \in U\)) and regular at every other prime divisor of \(U\), hence in
  particular a regular function. The cokernel is the skyscraper \(k(P)\) by
  \cref{lem:cokernel_sheafOf_single_add_iso_skyscraper}, via the residue
  (evaluation) map \(\mathcal O_C \twoheadrightarrow k(P)\) whose kernel is
  exactly the maximal-ideal sheaf \(\mathcal O_C(-[P])\). The sequence is
  therefore exact.
\end{proof}

% ════════════════════════════════════════════════════════════════════════
% The inductive-step SES is presented by the cokernel route rather than
% by the structure-sheaf tensor decomposition of Hartshorne~IV.1.3.
% The category $\mathbf{Sh}(C, \mathbf{Mod}_{\bar k})$ has $\bar k$-linear
% sheaves as objects (not $\mathcal O_C$-module objects), so the tensor
% $\otimes_{\mathcal O_C}$ is not its monoidal product and the classical
% ``tensor the ideal-sheaf sequence by $\mathcal O_C(D+[P])$'' argument
% does not apply.  The sequence is presented intrinsically: the subsheaf
% inclusion is a monomorphism, and the SES is the canonical
% mono--cokernel short exact sequence, with the cokernel identified as the
% skyscraper $k(P)$.
% ════════════════════════════════════════════════════════════════════════

\begin{definition}[Short complex and short-exactness]
  \label{def:shortComplex_mathlib}
  \lean{CategoryTheory.ShortComplex.ShortExact}
  \mathlibok
  \textit{Provided by Mathlib.}
  In a category with zero morphisms, a \emph{short complex} is a pair of
  composable morphisms \(X_1 \xrightarrow{f} X_2 \xrightarrow{g} X_3\) with
  \(g \circ f = 0\). In an abelian category it is \emph{short exact} when
  \(f\) is a monomorphism, \(g\) is an epimorphism, and the complex is exact
  at \(X_2\). This is Mathlib's \(\texttt{CategoryTheory.ShortComplex}\)
  together with its \(\texttt{ShortExact}\) predicate; it is the shape in
  which the project records the sequence
  \(0 \to \mathcal O_C(D) \to \mathcal O_C(D+[P]) \to k(P) \to 0\).
\end{definition}

\begin{definition}[Cokernel and its projection]
  \label{def:cokernel_mathlib}
  \lean{CategoryTheory.Limits.cokernel}
  \mathlibok
  \textit{Provided by Mathlib.}
  For a morphism \(f : X \to Y\) in a category with zero morphisms and the
  relevant colimit, the \emph{cokernel} \(\operatorname{coker} f\) carries a
  canonical projection \(\pi_f : Y \to \operatorname{coker} f\), characterised
  by \(\pi_f \circ f = 0\) and the universal property that any morphism out of
  \(Y\) killing \(f\) factors uniquely through \(\pi_f\). The abelian category
  \(\mathbf{Sh}(C, \mathbf{Mod}_{\bar k})\) has all cokernels.
\end{definition}

\begin{definition}[Skyscraper sheaf]
  \label{def:skyscraperSheaf_mathlib}
  \lean{skyscraperSheaf}
  \mathlibok
  \textit{Provided by Mathlib.}
  For a point \(p_0\) of a topological space \(X\) and an object \(A\) of a
  category with a terminal object, the skyscraper sheaf
  \(\texttt{skyscraperSheaf}\,p_0\,A\) takes the value \(A\) on every open
  containing \(p_0\) and the terminal object on every open avoiding \(p_0\).
  The third term of the inductive-step sequence is
  \(k(P) := \texttt{skyscraperSheaf}\,P.\mathrm{point}\,\bar k\) in
  \(\mathbf{Sh}(C, \mathbf{Mod}_{\bar k})\): the residue field \(\bar k\)
  placed at the closed point \(P\).
\end{definition}

\begin{lemma}

\leanok
[Adding \([P]\) relaxes the carrier order bound]
  \label{lem:carrierSet_le_add_single}
  \lean{AlgebraicGeometry.Scheme.WeilDivisor.sheafOf.carrierSet_le_add_single}
  \uses{def:sheafOf, def:order_at_point, def:prime_divisor}
  Let \(C\) be a smooth proper geometrically irreducible curve over
  \(\bar k\), let \(D \in \mathrm{Div}(C)\) and let \(P\) be a prime divisor
  of \(C\) (\cref{def:prime_divisor}). Over every open \(U \subseteq C\) the
  carrier set of \(\mathcal O_C(D)\) is contained in that of
  \(\mathcal O_C([P]+D)\):
  \[
    \Gamma\bigl(U, \mathcal O_C(D)\bigr) \;\subseteq\;
    \Gamma\bigl(U, \mathcal O_C([P]+D)\bigr)
  \]
  as subsets of \(K(C)\).
\end{lemma}

\begin{proof}
  \uses{def:order_at_point}
  The zero section lies in both carriers. A nonzero
  \(f \in \Gamma(U, \mathcal O_C(D))\) satisfies \(-(D\,Q) \leq \ord_Q(f)\)
  for every prime divisor \(Q \in U\). The coefficient of \(P\) in
  \([P]+D\) exceeds that in \(D\) by one, while every other coefficient is
  unchanged; hence \(-(([P]+D)\,Q) \leq -(D\,Q) \leq \ord_Q(f)\) at \(Q = P\)
  and the bounds are unchanged at \(Q \neq P\). Thus \(f\) satisfies the
  (weaker) order constraints defining \(\mathcal O_C([P]+D)\), and the
  inclusion of carrier sets follows.
\end{proof}

\begin{lemma}

\leanok
[Carrier-submodule inclusion in the divisor variable]
  \label{lem:carrierSubmoduleSheaf_le_add_single}
  \lean{AlgebraicGeometry.Scheme.WeilDivisor.sheafOf.carrierSubmoduleSheaf_le_add_single}
  \uses{lem:carrierSet_le_add_single, def:sheafOf}
  With \(C, D, P\) as above, on every open \(U\) the section
  \(\bar k\)-submodule of \(\mathcal O_C(D)\) is contained in that of
  \(\mathcal O_C([P]+D)\), as \(\bar k\)-submodules of \(K(C)\).
\end{lemma}

\begin{proof}
  \uses{lem:carrierSet_le_add_single}
  The defining submodule of \(\mathcal O_C(D)\) over \(U\) is the carrier set
  intersected with the bot-trivialisation constraint (which forces the
  section to vanish on the empty open). The carrier-set inclusion is
  \cref{lem:carrierSet_le_add_single}, and the bot-trivialisation factor
  depends on \(U\) alone, not on the divisor, so it coincides on both sides.
  Intersecting a smaller carrier with the same factor yields the asserted
  submodule inclusion.
\end{proof}

\begin{definition}

\leanok
[Inclusion morphism of carrier presheaves]
  \label{lem:carrierPresheaf_le_hom}
  \lean{AlgebraicGeometry.Scheme.WeilDivisor.sheafOf.carrierPresheaf_le_hom}
  \uses{lem:carrierSubmoduleSheaf_le_add_single, def:sheafOf}
  With \(C, D, P\) as above, the section-wise submodule inclusions of
  \cref{lem:carrierSubmoduleSheaf_le_add_single} assemble into a morphism of
  carrier presheaves
  \[
    \mathcal O_C(D) \;\longrightarrow\; \mathcal O_C([P]+D),
  \]
  natural in the open \(U\): the presheaf-level inclusion
  \(\mathcal O_C(D) \hookrightarrow \mathcal O_C([P]+D)\) of subpresheaves of
  the constant sheaf \(\mathcal K_C\).
\end{definition}

\begin{proof}
  \uses{lem:carrierSubmoduleSheaf_le_add_single}
  For each open \(U\) take the \(\bar k\)-linear submodule inclusion of
  \cref{lem:carrierSubmoduleSheaf_le_add_single}. Naturality is the
  commutation of these inclusions with restriction. The restriction maps on
  both carrier presheaves are the identity on \(K(C)\), so for
  \(V \subseteq U\) with \(V \neq \varnothing\) the naturality square
  commutes on the nose; at \(V = \varnothing\) the codomain section module is
  the zero module, a subsingleton, so the square commutes trivially. The
  section-wise inclusions therefore constitute a presheaf morphism.
\end{proof}

\begin{definition}

\leanok
[The carrier sheaf of \(\mathcal O_C(D)\)]
  \label{def:sheafOf_carrierSheaf}
  \lean{AlgebraicGeometry.Scheme.WeilDivisor.sheafOf.carrierSheaf}
  \uses{def:sheafOf_carrierPresheaf, lem:sheafOf_carrierPresheaf_isSheaf}
  For a Weil divisor \(D : C.\text{left.PrimeDivisor} \to_0 \mathbb Z\), the
  \emph{carrier sheaf} \(\texttt{carrierSheaf}_D\) is the presheaf
  \cref{def:sheafOf_carrierPresheaf} equipped with the sheaf property
  \cref{lem:sheafOf_carrierPresheaf_isSheaf}, regarded as an object of
  \(\mathbf{Sh}(C, \mathbf{Mod}_{\bar k})\). The lemmas below establish
  the inductive-step short exact sequence at the carrier-sheaf level.
\end{definition}

\begin{definition}

\leanok
[Sheaf-level inclusion \(\mathcal O_C(D) \hookrightarrow \mathcal O_C([P]+D)\)]
  \label{def:sheafOf_carrierSheafHom_le_add_single}
  \lean{AlgebraicGeometry.Scheme.WeilDivisor.sheafOf.carrierSheafHom_le_add_single}
  \uses{def:sheafOf_carrierSheaf, lem:carrierPresheaf_le_hom}
  With \(C, D\) as above and \(P\) a prime divisor of \(C\)
  (\cref{def:prime_divisor}), the presheaf inclusion
  \cref{lem:carrierPresheaf_le_hom} promotes to a morphism of carrier sheaves
  \[
    \texttt{carrierSheaf}_D \;\longrightarrow\;
    \texttt{carrierSheaf}_{[P]+D}
  \]
  in \(\mathbf{Sh}(C, \mathbf{Mod}_{\bar k})\): a morphism of sheaves is
  precisely a morphism of underlying presheaves, so the carrier-presheaf
  inclusion \(\mathcal O_C(D) \hookrightarrow \mathcal O_C([P]+D)\) is already a
  sheaf morphism between the carrier sheaves.
\end{definition}

\begin{lemma}

\leanok
[The inclusion \(\mathcal O_C(D) \hookrightarrow \mathcal O_C([P]+D)\) is a
monomorphism]
  \label{lem:sheafOf_mono_single_add}
  \lean{AlgebraicGeometry.Scheme.WeilDivisor.sheafOf.carrierSheafHom_le_add_single_mono}
  \uses{def:sheafOf_carrierSheafHom_le_add_single,
        lem:carrierSubmoduleSheaf_le_add_single, def:prime_divisor}
  Let \(C\) be a smooth proper geometrically irreducible curve over
  \(\bar k\), let \(D \in \mathrm{Div}(C)\) and let \(P\) be a closed point of
  \(C\) (a prime divisor, \cref{def:prime_divisor}). The sheaf-level inclusion
  of \cref{def:sheafOf_carrierSheafHom_le_add_single}
  \[
    f \;\colon\; \mathcal O_C(D) \;\longrightarrow\; \mathcal O_C([P]+D)
  \]
  in \(\mathbf{Sh}(C, \mathbf{Mod}_{\bar k})\) is a \emph{monomorphism}.
\end{lemma}

\begin{proof}
  \uses{def:sheafOf_carrierSheafHom_le_add_single,
        lem:carrierSubmoduleSheaf_le_add_single}
  Over every open \(U\) the morphism \(f\) is the \(\bar k\)-linear submodule
  inclusion of \cref{lem:carrierSubmoduleSheaf_le_add_single}, hence injective
  on sections; so each component of \(f\) is a monomorphism in
  \(\mathbf{Mod}_{\bar k}\). A natural transformation that is component-wise a
  monomorphism is a monomorphism of presheaves, and a presheaf monomorphism
  between sheaves is a monomorphism in the sheaf category
  \(\mathbf{Sh}(C, \mathbf{Mod}_{\bar k})\) (monomorphisms in this abelian
  sheaf category are detected section-wise). Therefore \(f\) is a monomorphism.
\end{proof}

\begin{lemma}\leanok
[Away from \(P\): the carrier set is unchanged by \([P]\)]
  \label{lem:carrierSet_add_single_eq_of_not_mem}
  \lean{AlgebraicGeometry.Scheme.WeilDivisor.sheafOf.carrierSet_add_single_eq_of_not_mem}
  \uses{def:sheafOf_carrierSet}
  Let \(C, D\) be as above and let \(P\) be a prime divisor of \(C\). For every
  open \(U \subseteq C\) with \(P.\mathrm{point} \notin U\) the carrier sets of
  \(\mathcal O_C([P]+D)\) and \(\mathcal O_C(D)\) over \(U\) coincide:
  \[
    \mathrm{carrierSet}_{[P]+D}(U) \;=\; \mathrm{carrierSet}_D(U).
  \]
\end{lemma}

\begin{proof}
  \uses{def:sheafOf_carrierSet}
  The divisors \([P]+D\) and \(D\) differ only in their coefficient at the prime
  \(P\). For every prime divisor \(Q\) with \(Q.\mathrm{point} \in U\) we have
  \(Q \neq P\) (else \(P.\mathrm{point} = Q.\mathrm{point} \in U\) by injectivity
  of \(Q \mapsto Q.\mathrm{point}\), contradicting \(P.\mathrm{point} \notin U\)), so \(([P]+D)\,Q = D\,Q\) and the
  order constraints defining the two carrier sets are identical over \(U\).
\end{proof}

\begin{lemma}\leanok
[Away from \(P\): the carrier submodule is unchanged]
  \label{lem:carrierSubmodule_add_single_eq_of_not_mem}
  \lean{AlgebraicGeometry.Scheme.WeilDivisor.sheafOf.carrierSubmodule_add_single_eq_of_not_mem}
  \uses{lem:carrierSet_add_single_eq_of_not_mem, def:sheafOf_carrierSubmodule}
  With \(C, D, P\) as above and \(U \subseteq C\) open with
  \(P.\mathrm{point} \notin U\), the carrier \(\bar k\)-submodules agree:
  \(\mathrm{carrierSubmodule}_{[P]+D}(U) = \mathrm{carrierSubmodule}_D(U)\).
\end{lemma}

\begin{proof}
  \uses{lem:carrierSet_add_single_eq_of_not_mem}
  The carrier submodule has underlying set the carrier set, so the equality of
  carrier sets \cref{lem:carrierSet_add_single_eq_of_not_mem} promotes to an
  equality of submodules by extensionality on the underlying sets.
\end{proof}

\begin{lemma}

\leanok
[Away from \(P\): the sheaf-corrected carrier submodule is unchanged]
  \label{lem:carrierSubmoduleSheaf_add_single_eq_of_not_mem}
  \lean{AlgebraicGeometry.Scheme.WeilDivisor.sheafOf.carrierSubmoduleSheaf_add_single_eq_of_not_mem}
  \uses{lem:carrierSubmodule_add_single_eq_of_not_mem,
        def:sheafOf_carrierSubmoduleSheaf}
  With \(C, D, P\) as above and \(U \subseteq C\) open with
  \(P.\mathrm{point} \notin U\), the sheaf-corrected carrier submodules agree:
  \(\mathrm{carrierSubmoduleSheaf}_{[P]+D}(U) =
  \mathrm{carrierSubmoduleSheaf}_D(U)\).
\end{lemma}

\begin{proof}
  \uses{lem:carrierSubmodule_add_single_eq_of_not_mem}
  The sheaf-corrected carrier submodule is the meet of the carrier submodule
  with the bot-trivialisation factor \(\mathrm{trivAtBot}(U)\), and the latter
  depends on \(U\) alone, not on the divisor. Rewriting the carrier-submodule
  factor by \cref{lem:carrierSubmodule_add_single_eq_of_not_mem} gives the
  asserted equality.
\end{proof}

\begin{lemma}\leanok
[Away from \(P\): the inclusion is an isomorphism on sections]
  \label{lem:carrierPresheaf_le_hom_app_isIso_of_not_mem}
  \lean{AlgebraicGeometry.Scheme.WeilDivisor.sheafOf.carrierPresheaf_le_hom_app_isIso_of_not_mem}
  \uses{lem:carrierSubmoduleSheaf_add_single_eq_of_not_mem,
        def:sheafOf_carrierSheafHom_le_add_single}
  With \(C, D, P\) as above and \(U \subseteq C\) open with
  \(P.\mathrm{point} \notin U\), the component over \(U\) of the inclusion
  morphism \(\mathcal O_C(D) \hookrightarrow \mathcal O_C([P]+D)\)
  (\cref{def:sheafOf_carrierSheafHom_le_add_single}) is an isomorphism of
  \(\bar k\)-modules.
\end{lemma}

\begin{proof}
  \uses{lem:carrierSubmoduleSheaf_add_single_eq_of_not_mem}
  By \cref{lem:carrierSubmoduleSheaf_add_single_eq_of_not_mem} the source and
  target section submodules coincide over \(U\), so the
  inclusion-of-equal-submodules component is the identity map, in particular a
  bijection, hence an isomorphism in \(\mathbf{Mod}_{\bar k}\).
\end{proof}

\begin{definition}

\leanok
[Order submodule at \(P\): the fractional ideal \(\pi_P^{-n}\mathcal O_{C,P}\)]
  \label{def:sheafOf_orderAtPSubmodule}
  \lean{AlgebraicGeometry.Scheme.WeilDivisor.sheafOf.orderAtPSubmodule}
  \uses{def:order_at_point, lem:stalk_isDVR_of_smooth}
  Let \(D : C.\mathrm{left.PrimeDivisor} \to_0 \mathbb Z\) and let \(P\) be a
  prime divisor of \(C\), with \(n = D\,P\). The \emph{order submodule at \(P\)}
  is the \(\bar k\)-submodule of \(K(C)\)
  \[
    \mathrm{orderAtP}_D(P) \;:=\;
    \bigl\{\, g \in K(C) \;\big|\; g = 0 \,\lor\, -\,n \le \ord_P(g) \,\bigr\}
    \;\subseteq\; K(C),
  \]
  cut out by the \emph{single} order constraint at \(P\) (with the separate
  \(g = 0\) clause, as in \cref{def:sheafOf_carrierSet}). In the discrete
  valuation ring \(\mathcal O_{C,P}\) (\cref{lem:stalk_isDVR_of_smooth}) with
  uniformiser \(\pi_P\) this submodule is exactly the fractional ideal
  \(\pi_P^{-n}\mathcal O_{C,P}\): a nonzero \(g\) has \(\ord_P(g) \ge -n\) iff
  \(\pi_P^{\,n} g \in \mathcal O_{C,P}\). It is the closure of the carrier
  conditions to the lone prime \(P\), and is the target of the carrier-stalk
  identification \cref{lem:carrierSheaf_stalk_eq_fractionalIdeal}.
\end{definition}

\begin{lemma}

\leanok
[Stalk of an inclusion-restriction presheaf is the directed supremum of its
sections]
  \label{lem:carrierSheaf_stalk_iso_iSup}
  \lean{AlgebraicGeometry.Scheme.WeilDivisor.sheafOf.carrierSheaf_stalk_iso_iSup}
  \uses{def:sheafOf_carrierSheaf, def:sheafOf_carrierSubmoduleSheaf,
        lem:sheafOf_carrierSubmoduleSheaf_le}
  With \(C, D\) as above and \(P\) a prime divisor, the stalk of the carrier
  sheaf at \(P.\mathrm{point}\) is the directed supremum, taken inside \(K(C)\),
  of the section submodules over the open neighbourhoods of \(P.\mathrm{point}\):
  \[
    (\texttt{carrierSheaf}_D)_{P.\mathrm{point}} \;\cong\;
    \bigsqcup_{U \ni P.\mathrm{point}}
        \mathrm{carrierSubmoduleSheaf}_D(U)
  \]
  as \(\bar k\)-modules, where the supremum ranges over open
  neighbourhoods \(U \ni P.\mathrm{point}\) and \(\bigsqcup\) is the lattice
  supremum of \(\bar k\)-submodules of \(K(C)\).
\end{lemma}

\begin{proof}
  \uses{def:sheafOf_carrierSubmoduleSheaf, lem:sheafOf_carrierSubmoduleSheaf_le,
        lem:carrierPresheaf_map_coe}
  Write \(\Gamma(U) := \mathrm{carrierSubmoduleSheaf}_D(U) \subseteq K(C)\) and
  \(S := \bigsqcup_{U \ni P.\mathrm{point}} \Gamma(U)\). The family
  \(U \mapsto \Gamma(U)\) over \(\mathrm{OpenNhds}(P.\mathrm{point})\) is
  \emph{directed} under \(\le\): the index is nonempty (\(\top \ni
  P.\mathrm{point}\)) and for \(U, V \ni P.\mathrm{point}\) the open \(U \sqcap V
  \ni P.\mathrm{point}\) satisfies \(\Gamma(U), \Gamma(V) \le \Gamma(U \sqcap
  V)\) by monotonicity (\cref{lem:sheafOf_carrierSubmoduleSheaf_le}; note
  \(U \sqcap V \ne \emptyset\) as it contains \(P.\mathrm{point}\)). Since the
  family is directed, the supremum \(S\) equals the genuine union
  \(\bigcup_{U \ni P.\mathrm{point}} \Gamma(U)\) inside \(K(C)\).

  Every restriction map of \(\texttt{carrierSheaf}_D\) between nonempty opens is
  the inclusion \(\Gamma(U) \hookrightarrow \Gamma(V)\) of submodules of \(K(C)\)
  (\cref{def:sheafOf_carrierSubmoduleSheaf}), i.e.\ the identity on the
  underlying elements of \(K(C)\). We build the comparison map
  \(\varphi : (\texttt{carrierSheaf}_D)_{P.\mathrm{point}} \to S\) via the
  universal property of the stalk: each section leg \(\Gamma(U) \hookrightarrow
  S\) is the canonical inclusion into the supremum, and these maps are compatible
  with the (inclusion) restriction maps, so they factor through the colimit
  defining the stalk; on the germ of \(s \in \Gamma(U)\) at \(P.\mathrm{point}\),
  \(\varphi\) sends it to \(s\) viewed in \(S\). Then:
  \begin{itemize}
  \item \(\varphi\) is \emph{injective}: post-composing with the inclusion
    \(S \hookrightarrow K(C)\) sends the germ of \(s\) over \(U\) to the
    underlying element of \(s\) in \(K(C)\); since every restriction map is an
    injective inclusion, two germs with equal underlying element agree on a
    common smaller neighbourhood and are therefore equal by sheaf separatedness.
  \item \(\varphi\) is \emph{surjective}: every element of \(S\) lies in some
    \(\Gamma(U)\) with \(U \ni P.\mathrm{point}\) by the directed-union
    description above, and is thus the image under \(\varphi\) of the
    corresponding germ.
  \end{itemize}
  A bijective \(\bar k\)-linear map of modules is an isomorphism, giving the
  stated isomorphism.
\end{proof}

\begin{lemma}

\leanok
[Finitely many prime divisors violate the order constraints on an affine open]
  \label{lem:finite_orderConstraintFail_affine}
  \lean{AlgebraicGeometry.Scheme.WeilDivisor.sheafOf.finite_orderConstraintFail_affine}
  % NOTE: The proof uses global finiteness of the principal-divisor support
  % (available on the proper curve via compactness), which is strictly stronger
  % than the affine-finiteness route and handles \(g = 0\) internally --- so
  % the statement carries no nonzero hypothesis on \(g\).
  \uses{def:order_at_point, def:prime_divisor, def:principal_divisor,
        def:codim1_cycles}
  Let \(W \subseteq C\) be an affine open and let \(g \in K(C)\) be a rational
  function. Then the set of prime divisors
  \[
    \bigl\{\, Q \;\big|\; Q.\mathrm{point} \in W \,\land\,
      \neg\,(-\,D\,Q \le \ord_Q(g)) \,\bigr\}
  \]
  is finite.
\end{lemma}

\begin{proof}
  \uses{def:principal_divisor}
  If a prime divisor \(Q\) with \(Q.\mathrm{point} \in W\) violates the
  constraint, then \(\ord_Q(g) < -\,D\,Q\); in particular it is impossible to
  have both \(\ord_Q(g) = 0\) and \(D\,Q = 0\) (which would give \(0 \le 0\)).
  Hence every such \(Q\) lies in the union of the two sets
  \[
    \{\, Q : \ord_Q(g) \ne 0 \,\}
    \quad\text{and}\quad
    \{\, Q : D\,Q \ne 0 \,\}.
  \]
  The second set is finite because \(D\) has finite support by definition of a
  Weil divisor. The first set is finite because it is the support of the
  principal divisor \(\operatorname{div}(g)\) (\cref{def:principal_divisor}):
  on the proper curve \(C\) the curve's underlying space is compact, so
  \(\operatorname{div}(g)\) is a well-defined finitely-supported cycle whose
  support is exactly \(\{\, Q : \ord_Q(g) \ne 0 \,\}\) (when \(g = 0\) this set
  is empty by the convention \(\ord_Q(0) = 0\)). Being the support of a
  \(\mathrm{Finsupp}\), it is finite outright — no affine restriction is
  needed. A finite union of finite sets is finite.
\end{proof}

\begin{lemma}

\leanok
[Order constraint at \(P\) is realised on a small open]
  \label{lem:exists_open_mem_carrierSubmoduleSheaf}
  \lean{AlgebraicGeometry.Scheme.WeilDivisor.sheafOf.exists_open_mem_carrierSubmoduleSheaf}
  \uses{def:sheafOf_orderAtPSubmodule, def:sheafOf_carrierSubmoduleSheaf,
        lem:finite_orderConstraintFail_affine, def:prime_divisor}
  Let \(g \in \mathrm{orderAtP}_D(P)\) (\cref{def:sheafOf_orderAtPSubmodule}).
  Then there is an open \(U \ni P.\mathrm{point}\) with
  \(g \in \mathrm{carrierSubmoduleSheaf}_D(U)\).
\end{lemma}

\begin{proof}
  \uses{lem:finite_orderConstraintFail_affine}
  If \(g = 0\) then \(g \in \mathrm{carrierSubmoduleSheaf}_D(U)\) for every
  \(U \ni P.\mathrm{point}\) (the \(f = 0\) clause), so take \(U = \top\).
  Assume \(g \ne 0\). Choose an affine open \(W \ni P.\mathrm{point}\) (every
  point of a scheme has an affine open neighbourhood). By
  \cref{lem:finite_orderConstraintFail_affine} the set
  \[
    B \;=\; \{\, Q \mid Q.\mathrm{point} \in W \,\land\,
      \neg\,(-\,D\,Q \le \ord_Q(g)) \,\}
  \]
  of \emph{bad} prime divisors inside \(W\) is finite. Each \(Q \in B\) is a
  codimension-one point of \(C\), hence its closure \(\overline{\{Q.\mathrm
  {point}\}}\) is a proper closed set; moreover \(P \notin B\), because
  \(g \in \mathrm{orderAtP}_D(P)\) and \(g \ne 0\) give exactly
  \(-\,D\,P \le \ord_P(g)\). Therefore
  \[
    U \;:=\; W \;\setminus\; \bigcup_{Q \in B}
      \overline{\{Q.\mathrm{point}\}}
  \]
  is open (a finite union of closed sets removed from the open \(W\)) and
  contains \(P.\mathrm{point}\) (it lies in \(W\), and \(P.\mathrm{point}\) is
  not in any \(\overline{\{Q.\mathrm{point}\}}\) for \(Q \in B\): a prime
  divisor's generic point lies in another's closure only if they coincide, and
  \(P \notin B\)). For every prime divisor \(Q'\) with \(Q'.\mathrm{point} \in
  U\) we have \(Q'.\mathrm{point} \in W\) and \(Q' \notin B\), so
  \(-\,D\,Q' \le \ord_{Q'}(g)\). Hence \(g\) satisfies all order constraints
  over \(U\), i.e.\ \(g \in \mathrm{carrierSubmoduleSheaf}_D(U)\) (and \(U \ne
  \emptyset\) since \(P.\mathrm{point} \in U\), so the \(\mathrm{trivAtBot}\)
  factor is \(K(C)\)).
\end{proof}

\begin{lemma}

\leanok
[The directed supremum of carrier sections at \(P\) is the order submodule]
  \label{lem:iSup_carrierSubmoduleSheaf_eq_orderAtPSubmodule}
  \lean{AlgebraicGeometry.Scheme.WeilDivisor.sheafOf.iSup_carrierSubmoduleSheaf_eq_orderAtPSubmodule}
  \uses{def:sheafOf_carrierSubmoduleSheaf, def:sheafOf_orderAtPSubmodule,
        lem:sheafOf_carrierSubmoduleSheaf_le,
        lem:exists_open_mem_carrierSubmoduleSheaf, def:prime_divisor}
  With \(C, D, P\) as above,
  \[
    \bigsqcup_{U \ni P.\mathrm{point}}
      \mathrm{carrierSubmoduleSheaf}_D(U) \;=\; \mathrm{orderAtP}_D(P)
  \]
  as \(\bar k\)-submodules of \(K(C)\).
\end{lemma}

\begin{proof}
  \uses{lem:exists_open_mem_carrierSubmoduleSheaf,
        lem:sheafOf_carrierSubmoduleSheaf_le}
  We prove the two inclusions; throughout, the family is directed (as in
  \cref{lem:carrierSheaf_stalk_iso_iSup}), so membership in the supremum is
  equivalent to membership in some \(\mathrm{carrierSubmoduleSheaf}_D(U)\) with
  \(U \ni P.\mathrm{point}\).

  \emph{(\(\le\))} It suffices to show each leg
  \(\mathrm{carrierSubmoduleSheaf}_D(U) \le \mathrm{orderAtP}_D(P)\) for \(U \ni
  P.\mathrm{point}\). Let \(g \in \mathrm{carrierSubmoduleSheaf}_D(U)\). If
  \(g = 0\) it lies in \(\mathrm{orderAtP}_D(P)\). Otherwise, since
  \(P.\mathrm{point} \in U\), the order constraint at \(Q = P\) is among those
  enforced over \(U\), giving \(-\,D\,P \le \ord_P(g)\); thus
  \(g \in \mathrm{orderAtP}_D(P)\).

  \emph{(\(\ge\))} Let \(g \in \mathrm{orderAtP}_D(P)\). By
  \cref{lem:exists_open_mem_carrierSubmoduleSheaf} there is an open \(U \ni
  P.\mathrm{point}\) with \(g \in \mathrm{carrierSubmoduleSheaf}_D(U)\), whence
  \(g\) lies in the directed supremum.
\end{proof}

\begin{lemma}\leanok
[Stalk of the carrier sheaf is a fractional ideal of the DVR]
  \label{lem:carrierSheaf_stalk_eq_fractionalIdeal}
  \lean{AlgebraicGeometry.Scheme.WeilDivisor.sheafOf.carrierSheaf_stalk_eq}
  \uses{def:sheafOf_carrierSheaf, lem:stalk_isDVR_of_smooth, def:order_at_point,
        def:sheafOf_orderAtPSubmodule, lem:carrierSheaf_stalk_iso_iSup,
        lem:iSup_carrierSubmoduleSheaf_eq_orderAtPSubmodule}
  Let \(C\) be a smooth proper geometrically irreducible curve over \(\bar k\),
  let \(D \in \mathrm{Div}(C)\), and let \(P\) be a prime divisor of \(C\) with
  order \(n = D\,P\) (\cref{def:order_at_point}). The stalk of the carrier sheaf
  \(\texttt{carrierSheaf}_D\) (\cref{def:sheafOf_carrierSheaf}) at the point
  \(P.\mathrm{point}\) is the fractional ideal
  \[
    (\texttt{carrierSheaf}_D)_{P.\mathrm{point}} \;=\;
    \pi_P^{-n}\,\mathcal O_{C,P} \;\subseteq\; K(C),
  \]
  where \(\pi_P\) is a uniformiser of the discrete valuation ring
  \(\mathcal O_{C,P}\) (\cref{lem:stalk_isDVR_of_smooth}).
\end{lemma}

\begin{proof}
  \uses{lem:carrierSheaf_stalk_iso_iSup,
        lem:iSup_carrierSubmoduleSheaf_eq_orderAtPSubmodule,
        def:sheafOf_orderAtPSubmodule, lem:stalk_isDVR_of_smooth}
  By \cref{lem:carrierSheaf_stalk_iso_iSup} the stalk of the carrier sheaf at
  \(P.\mathrm{point}\) is the directed supremum
  \(\bigsqcup_{U \ni P.\mathrm{point}} \mathrm{carrierSubmoduleSheaf}_D(U)\),
  taken inside \(K(C)\) (every restriction map is the identity inclusion on
  \(K(C)\)). By
  \cref{lem:iSup_carrierSubmoduleSheaf_eq_orderAtPSubmodule} this supremum equals
  the order submodule \(\mathrm{orderAtP}_D(P)\) of
  \cref{def:sheafOf_orderAtPSubmodule}: the \(\le\) inclusion uses the constraint
  at \(Q = P\) on each leg, while the \(\ge\) inclusion realises any
  \(g \in \mathrm{orderAtP}_D(P)\) on an open neighbourhood of \(P.\mathrm{point}\)
  obtained by deleting the \emph{finitely many} prime divisors at which \(g\) or
  \(D\) is supported (its poles and the support of \(D\)), so that no other order
  constraint survives. Rewriting the supremum by this equality through the iso of
  \cref{lem:carrierSheaf_stalk_iso_iSup} identifies the stalk with
  \(\mathrm{orderAtP}_D(P)\). Finally, in the discrete valuation ring
  \(\mathcal O_{C,P}\) (\cref{lem:stalk_isDVR_of_smooth}) with uniformiser
  \(\pi_P\), the order submodule \(\mathrm{orderAtP}_D(P) = \{g \mid g = 0 \lor
  -n \le \ord_P(g)\}\) is precisely the fractional ideal
  \(\pi_P^{-n}\mathcal O_{C,P} \subseteq K(C)\), which is the stated form.

  The stalk is identified \emph{up to isomorphism} with
  \(\mathrm{orderAtP}_D(P)\) (\cref{def:sheafOf_orderAtPSubmodule}); the
  displayed equality \(=\) in the statement is therefore an isomorphism
  \(\cong\). The fractional-ideal description \(\pi_P^{-n}\mathcal O_{C,P}\)
  is the intrinsic characterisation of \(\mathrm{orderAtP}_D(P)\) given in
  \cref{def:sheafOf_orderAtPSubmodule}, not a separate equality.
\end{proof}

\begin{lemma}\leanok
[At \(P\): the cokernel stalk is \(\bar k\)]
  \label{lem:cokernel_sheafOf_single_add_stalkAtP_iso_kbar}
  \lean{AlgebraicGeometry.Scheme.WeilDivisor.sheafOf.cokernel_stalk_at_iso_kbar}
  \uses{lem:carrierSheaf_stalk_eq_fractionalIdeal, lem:stalk_isDVR_of_smooth,
        lem:codimOne_point_residueField_eq_kbar, def:order_at_point}
  With \(C, D, P\) as above and \(n = D\,P\), at the point \(P.\mathrm{point}\)
  the stalk of the inclusion \(\mathcal O_C(D) \hookrightarrow \mathcal
  O_C([P]+D)\) is the inclusion of fractional ideals
  \(\pi_P^{-n}\mathcal O_{C,P} \subseteq \pi_P^{-(n+1)}\mathcal O_{C,P}\) in
  \(K(C)\). Multiplication by \(\pi_P^{\,n+1}\) is a \(\bar k\)-linear
  automorphism of \(K(C)\) carrying this inclusion isomorphically onto the
  maximal-ideal inclusion \(\mathfrak m_P = \pi_P\mathcal O_{C,P} \subseteq
  \mathcal O_{C,P}\). Hence the stalk of the cokernel at \(P.\mathrm{point}\) is
  \[
    \pi_P^{-(n+1)}\mathcal O_{C,P} \big/ \pi_P^{-n}\mathcal O_{C,P}
    \;\cong\; \mathcal O_{C,P}/\mathfrak m_P \;=\; \bar k,
  \]
  the residue field of \(P\), which equals \(\bar k\) because \(P\) is a
  codimension-one point (\cref{lem:codimOne_point_residueField_eq_kbar}).
\end{lemma}

\begin{proof}
  \uses{lem:carrierSheaf_stalk_eq_fractionalIdeal, lem:stalk_isDVR_of_smooth,
        lem:codimOne_point_residueField_eq_kbar}
  By \cref{lem:carrierSheaf_stalk_eq_fractionalIdeal} the two stalks are
  \(\pi_P^{-n}\mathcal O_{C,P}\) and \(\pi_P^{-(n+1)}\mathcal O_{C,P}\), and the
  stalk of the inclusion is the inclusion of these fractional ideals (the
  restriction maps being the identity on \(K(C)\)). In the discrete valuation
  ring \(\mathcal O_{C,P}\) (\cref{lem:stalk_isDVR_of_smooth}) multiplication by
  \(\pi_P^{\,n+1}\) is a \(\bar k\)-linear bijection of \(K(C)\) sending
  \(\pi_P^{-n}\mathcal O_{C,P}\) onto \(\pi_P\mathcal O_{C,P} = \mathfrak m_P\)
  and \(\pi_P^{-(n+1)}\mathcal O_{C,P}\) onto \(\mathcal O_{C,P}\). It therefore
  induces an isomorphism of quotient \(\bar k\)-modules
  \(\pi_P^{-(n+1)}\mathcal O_{C,P}/\pi_P^{-n}\mathcal O_{C,P} \cong \mathcal
  O_{C,P}/\mathfrak m_P\). The right-hand quotient is the residue field at \(P\),
  equal to \(\bar k\) by \cref{lem:codimOne_point_residueField_eq_kbar}.
\end{proof}

\begin{definition}\leanok
[Residue-evaluation comparison morphism]
  \label{def:cokernel_skyscraper_hom}
  \lean{AlgebraicGeometry.Scheme.WeilDivisor.sheafOf.cokernel_skyscraper_hom}
  \uses{def:sheafOf_carrierSheafHom_le_add_single, def:cokernel_mathlib,
        def:skyscraperSheaf_mathlib, def:order_at_point}
  With \(C, D, P\) as above and \(n = D\,P\), the \emph{residue-evaluation
  morphism} is the comparison morphism
  \[
    \operatorname{coker} f \;\longrightarrow\;
    \texttt{skyscraperSheaf}\,P.\mathrm{point}\,\bar k
  \]
  in \(\mathbf{Sh}(C, \mathbf{Mod}_{\bar k})\), where \(f\) is the sheaf-level
  inclusion of \cref{def:sheafOf_carrierSheafHom_le_add_single} and
  \(\operatorname{coker} f\) its cokernel (\cref{def:cokernel_mathlib}). On a
  small open \(U \ni P.\mathrm{point}\) it sends the class of a section
  \(g \in \mathcal O_C([P]+D)(U)\) to its leading Laurent coefficient at the
  order-\((n+1)\) pole at \(P\) --- the residue
  \((\pi_P^{\,n+1} g) \bmod \mathfrak m_P \in \bar k\) --- and it is the zero map
  on every open avoiding \(P.\mathrm{point}\). This evaluation is well-defined
  modulo \(\mathcal O_C(D)(U)\) and natural in \(U\), so it descends from the
  cokernel to the skyscraper.
\end{definition}

\begin{lemma}\leanok
[The residue-evaluation morphism is an isomorphism]
  \label{lem:cokernel_skyscraper_hom_isIso}
  \lean{AlgebraicGeometry.Scheme.WeilDivisor.sheafOf.cokernel_skyscraper_hom_isIso}
  \uses{def:cokernel_skyscraper_hom,
        lem:carrierPresheaf_le_hom_app_isIso_of_not_mem,
        lem:cokernel_sheafOf_single_add_stalkAtP_iso_kbar,
        def:skyscraperSheaf_mathlib}
  The residue-evaluation morphism \cref{def:cokernel_skyscraper_hom} induces an
  isomorphism on the stalk at every point of \(C\): the zero map at every point
  other than \(P.\mathrm{point}\) (both stalks vanish there), and the residue
  isomorphism \(\mathcal O_{C,P}/\mathfrak m_P = \bar k\) at
  \(P.\mathrm{point}\). A morphism of sheaves that is an isomorphism on every
  stalk is itself an isomorphism, so the residue-evaluation morphism is an
  isomorphism
  \(\operatorname{coker} f \cong
  \texttt{skyscraperSheaf}\,P.\mathrm{point}\,\bar k\).
\end{lemma}

\begin{proof}
  \uses{lem:carrierPresheaf_le_hom_app_isIso_of_not_mem,
        lem:cokernel_sheafOf_single_add_stalkAtP_iso_kbar}
  Stalks are checked pointwise. At a point \(Q \neq P.\mathrm{point}\): on a
  small open \(U \ni Q\) with \(P.\mathrm{point} \notin U\) the inclusion \(f\) is
  an isomorphism on sections by
  \cref{lem:carrierPresheaf_le_hom_app_isIso_of_not_mem}, so \(\operatorname{coker}
  f\) has zero stalk at \(Q\); the skyscraper likewise has zero (terminal) stalk
  away from \(P.\mathrm{point}\), and the morphism between zero modules is an
  isomorphism. At \(Q = P.\mathrm{point}\): the stalk of \(\operatorname{coker}
  f\) is \(\bar k\) by
  \cref{lem:cokernel_sheafOf_single_add_stalkAtP_iso_kbar}, and the
  residue-evaluation morphism realises precisely the residue isomorphism
  \(\mathcal O_{C,P}/\mathfrak m_P = \bar k\) onto the skyscraper stalk
  \(\bar k\). Being a stalkwise isomorphism, the morphism is an isomorphism of
  sheaves.
\end{proof}

\begin{lemma}

\leanok
[The cokernel of the inclusion is the skyscraper at \(P\)]
  \label{lem:cokernel_sheafOf_single_add_iso_skyscraper}
  \lean{AlgebraicGeometry.Scheme.WeilDivisor.sheafOf.cokernel_carrierSheafHom_iso_skyscraper}
  \uses{lem:sheafOf_mono_single_add, def:cokernel_mathlib,
        def:skyscraperSheaf_mathlib,
        def:cokernel_skyscraper_hom,
        lem:cokernel_sheafOf_single_add_stalkAtP_iso_kbar,
        lem:cokernel_skyscraper_hom_isIso}
  Let \(C, D, P\) be as above and let
  \(f : \mathcal O_C(D) \hookrightarrow \mathcal O_C([P]+D)\) be the
  monomorphism of \cref{lem:sheafOf_mono_single_add}. Then its cokernel in
  \(\mathbf{Sh}(C, \mathbf{Mod}_{\bar k})\) is isomorphic to the skyscraper
  sheaf at \(P\) with stalk \(\bar k\):
  \[
    \operatorname{coker} f \;\cong\;
    \texttt{skyscraperSheaf}\;P.\mathrm{point}\;\bar k.
  \]
\end{lemma}

\begin{proof}
  \uses{def:cokernel_skyscraper_hom,
        lem:cokernel_sheafOf_single_add_stalkAtP_iso_kbar,
        lem:cokernel_skyscraper_hom_isIso}
  Assemble the stalkwise analysis. The residue-evaluation comparison morphism
  \cref{def:cokernel_skyscraper_hom} from \(\operatorname{coker} f\) to the
  skyscraper \(\texttt{skyscraperSheaf}\,P.\mathrm{point}\,\bar k\) is an
  isomorphism on every stalk: it is the zero-to-zero map away from
  \(P.\mathrm{point}\), and at \(P.\mathrm{point}\) it realises the residue
  isomorphism \(\mathcal O_{C,P}/\mathfrak m_P = \bar k\) computed in
  \cref{lem:cokernel_sheafOf_single_add_stalkAtP_iso_kbar}. By
  \cref{lem:cokernel_skyscraper_hom_isIso} a stalkwise isomorphism of sheaves is
  an isomorphism, so \(\operatorname{coker} f \cong
  \texttt{skyscraperSheaf}\,P.\mathrm{point}\,\bar k = k(P)\).
\end{proof}

\begin{lemma}

\leanok
[\(\mathcal O_C(D)\) is its carrier sheaf for a nonzero divisor]
  \label{lem:sheafOf_eq_carrier_of_ne_zero}
  \lean{AlgebraicGeometry.Scheme.WeilDivisor.sheafOf_eq_carrier_of_ne_zero}
  \uses{def:sheafOf, def:sheafOf_carrierSheaf}
  Let \(C\) be a smooth proper geometrically irreducible curve over \(\bar k\)
  and let \(D \in \mathrm{Div}(C)\) with \(D \neq 0\). Then the invertible
  sheaf \(\mathcal O_C(D)\) of \cref{def:sheafOf} equals its carrier sheaf
  of \cref{def:sheafOf_carrierSheaf}:
  \[
    \mathcal O_C(D) = \texttt{carrierSheaf}_D.
  \]
  This equality transports the carrier-level inductive-step results to the
  level of \(\mathcal O_C(D)\).
\end{lemma}

\begin{proof}
  \uses{def:sheafOf, def:sheafOf_carrierSheaf}
  By \cref{def:sheafOf}, when \(D \neq 0\) the invertible sheaf is defined
  as \(\texttt{carrierSheaf}_D\), giving the asserted equality.
\end{proof}

\begin{theorem}

\leanok
[Short exact sequence for \(D \rightsquigarrow D + [P]\) at the carrier level]
  \label{thm:sheafOf_carrierSheaf_ses_single_add}
  \lean{AlgebraicGeometry.Scheme.WeilDivisor.sheafOf.carrierSheaf_ses_single_add}
  \uses{def:sheafOf_carrierSheaf,
        def:sheafOf_carrierSheafHom_le_add_single,
        lem:sheafOf_mono_single_add,
        lem:cokernel_sheafOf_single_add_iso_skyscraper,
        def:shortComplex_mathlib, def:cokernel_mathlib,
        def:skyscraperSheaf_mathlib}
  Let \(C, D, P\) be as above. The sheaf-level inclusion
  \cref{def:sheafOf_carrierSheafHom_le_add_single}
  \(f : \texttt{carrierSheaf}_D \to \texttt{carrierSheaf}_{[P]+D}\) fits into a
  short exact sequence
  \[
    0 \;\longrightarrow\; \texttt{carrierSheaf}_D
    \;\longrightarrow\; \texttt{carrierSheaf}_{[P]+D}
    \;\longrightarrow\; k(P)
    \;\longrightarrow\; 0
  \]
  in \(\mathbf{Sh}(C, \mathbf{Mod}_{\bar k})\), with \(k(P)\) the skyscraper
  sheaf at \(P\) with stalk \(\bar k\). This is the carrier-sheaf formulation
  of \cref{lem:sheafOf_ses_single_add}, using the identification
  \(\mathcal O_C(D) = \texttt{carrierSheaf}_D\) of
  \cref{lem:sheafOf_eq_carrier_of_ne_zero} (valid when \(D \neq 0\)).
\end{theorem}

\begin{proof}
  \uses{def:sheafOf_carrierSheafHom_le_add_single,
        lem:sheafOf_mono_single_add,
        lem:cokernel_sheafOf_single_add_iso_skyscraper,
        def:shortComplex_mathlib, def:cokernel_mathlib}
  Take \(f\) the sheaf-level inclusion of
  \cref{def:sheafOf_carrierSheafHom_le_add_single}; it is a monomorphism by
  \cref{lem:sheafOf_mono_single_add}. Form the short complex
  (\cref{def:shortComplex_mathlib})
  \(\texttt{carrierSheaf}_D \xrightarrow{f} \texttt{carrierSheaf}_{[P]+D}
  \xrightarrow{\pi_f} \operatorname{coker} f\) with \(\pi_f\) the canonical
  cokernel projection (\cref{def:cokernel_mathlib}). In the abelian category
  \(\mathbf{Sh}(C, \mathbf{Mod}_{\bar k})\) the sequence attached to a
  monomorphism \(f\) is short exact: \(f\) is mono, \(\pi_f\) is epi as a
  cokernel projection, and exactness at the middle term is the defining
  property of the cokernel of a monomorphism. Finally
  \cref{lem:cokernel_sheafOf_single_add_iso_skyscraper} identifies the third
  term \(\operatorname{coker} f \cong k(P)\), giving the asserted sequence.
\end{proof}

\begin{lemma}[Algebraic Hartogs: bounded valuations give integral elements]
  \label{lem:mem_integers_of_valuation_le_one}
  \lean{IsDedekindDomain.HeightOneSpectrum.mem_integers_of_valuation_le_one}
  \mathlibok
  \textit{Provided by Mathlib (\texttt{Mathlib.RingTheory.DedekindDomain.AdicValuation}).}
  Let \(R\) be a Dedekind domain with fraction field
  \(K = \operatorname{Frac}(R)\), and let \(x \in K\). If the
  \(\mathfrak p\)-adic valuation satisfies \(v_{\mathfrak p}(x) \le 1\) at every
  height-one prime \(\mathfrak p\) of \(R\), then \(x\) lies in the image of the
  structure map \(R \hookrightarrow K\), i.e.\
  \(x \in (\operatorname{algebraMap} R\,K).\mathrm{range}\). Equivalently, inside
  \(K\) one has \(R = \bigcap_{\operatorname{ht}\mathfrak p = 1} R_{\mathfrak p}\).
  This is the ring-theoretic core ("algebraic Hartogs") behind
  Hartshorne's Proposition~6.3A on a normal Noetherian domain.
\end{lemma}

\begin{lemma}[Sections of the zero carrier sheaf are the structure sheaf]
  \label{lem:carrierSheaf_zero_sections_eq_structureSheaf}
  % NOTE: `\lean{}` names the EXPECTED Lean target for the section-wise
  % identification step; it isolates the single real Mathlib gap (the
  % algebraic-Hartogs anchor \cref{lem:mem_integers_of_valuation_le_one})
  % from the sheaf-theoretic assembly in
  % \cref{lem:carrierSheaf_zero_iso_toModuleKSheaf}.
  \lean{AlgebraicGeometry.Scheme.WeilDivisor.sheafOf.carrierSheaf_zero_sections_eq_structureSheaf}
  % SOURCE: [Hartshorne], II.6, Proposition 6.3A, p.~132 (read from
  % references/hartshorne-algebraic-geometry.pdf, PDF page 132 = doc p.132,
  % transcribed character-by-character from the rendered scanned page image;
  % the PDF has no text layer).
  % SOURCE QUOTE: "Proposition 6.3A. Let \(A\) be an integrally closed
  % noetherian domain. Then
  % \(A = \bigcap_{\mathrm{ht}\,\mathfrak p = 1} A_{\mathfrak p}\)
  % where the intersection is taken over all prime ideals of height 1.
  % PROOF. Matsumura [2, Th. 38, p. 124]."
  \textit{Source: Hartshorne, II.6, Proposition 6.3A, p.~132.}
  \uses{lem:mem_integers_of_valuation_le_one, def:order_at_point,
        def:sheafOf_carrierSheaf}
  Let \(C\) be a smooth proper geometrically irreducible curve over \(\bar k\)
  and let \(U = \operatorname{Spec} A \subseteq C\) be an affine open with
  coordinate ring \(A\). Then the section module of the carrier sheaf of the
  zero divisor coincides with the regular functions on \(U\):
  \[
    \Gamma(U, \texttt{carrierSheaf}_0) \;=\; \Gamma(U, \struct C).
  \]
\end{lemma}

% SOURCE QUOTE PROOF: "Proposition 6.3A. Let \(A\) be an integrally closed
% noetherian domain. Then
% \(A = \bigcap_{\mathrm{ht}\,\mathfrak p = 1} A_{\mathfrak p}\)
% where the intersection is taken over all prime ideals of height 1.
% PROOF. Matsumura [2, Th. 38, p. 124]."
\begin{proof}
  \uses{lem:mem_integers_of_valuation_le_one, def:order_at_point,
        def:sheafOf_carrierSheaf}
  Over the affine open \(U = \operatorname{Spec} A\) the section module of
  \(\texttt{carrierSheaf}_0\) is the set of rational functions with no poles on
  \(U\):
  \[
    \Gamma(U, \texttt{carrierSheaf}_0) \;=\;
    \{\, g \in K(C) : \ord_Q(g) \geq 0 \text{ for every prime divisor }
    Q \in U \,\},
  \]
  the order constraint at the zero divisor (\cref{def:order_at_point}). A smooth
  curve is regular, hence normal, so the coordinate ring \(A\) is an integrally
  closed Noetherian domain; being one-dimensional it is a Dedekind domain, with
  fraction field \(K(C)\). On the smooth curve the codimension-one points of
  \(U\) are precisely the prime divisors meeting \(U\), and under this
  identification the \(\mathfrak p\)-adic valuation
  \(v_{\mathfrak p} = \texttt{HeightOneSpectrum.valuation}\) at a height-one
  prime \(\mathfrak p\) agrees with the divisorial order \(\ord_Q\) at the
  corresponding prime divisor \(Q\). Thus a rational function \(g\) has
  \(\ord_Q(g) \ge 0\) at every prime divisor of \(U\) precisely when
  \(v_{\mathfrak p}(g) \le 1\) at every height-one prime of \(A\); by the
  algebraic-Hartogs anchor \cref{lem:mem_integers_of_valuation_le_one} this
  holds exactly when \(g\) lies in \(A = \Gamma(U, \struct C)\). Hence the two
  section modules coincide.
\end{proof}

\begin{lemma}\leanok
[The carrier sheaf of the zero divisor is the structure sheaf]
  \label{lem:carrierSheaf_zero_iso_toModuleKSheaf}
  \lean{AlgebraicGeometry.Scheme.WeilDivisor.sheafOf.carrierSheaf_zero_iso_toModuleKSheaf}
  \textit{Source: Hartshorne, II.6, Proposition 6.3A, p.~132.}
  \uses{def:sheafOf_carrierSheaf, def:Scheme_toModuleKSheaf, def:order_at_point}
  Let \(C\) be a smooth proper geometrically irreducible curve over \(\bar k\).
  The carrier sheaf of the zero divisor is isomorphic, in
  \(\mathbf{Sh}(C, \mathbf{Mod}_{\bar k})\), to the structure sheaf of \(C\)
  regarded as a sheaf of \(\bar k\)-modules \(\toModuleKSheaf C\)
  (\cref{def:Scheme_toModuleKSheaf}):
  \[
    \texttt{carrierSheaf}_0 \;\cong\; \toModuleKSheaf C.
  \]
\end{lemma}

\begin{proof}
  \uses{lem:carrierSheaf_zero_sections_eq_structureSheaf,
        def:Scheme_toModuleKSheaf}
  By \cref{lem:carrierSheaf_zero_sections_eq_structureSheaf} the section module
  of \(\texttt{carrierSheaf}_0\) over each affine open
  \(U = \operatorname{Spec} A \subseteq C\) equals \(\Gamma(U, \struct C)\), the
  regular functions on \(U\); both are realised as \(\bar k\)-submodules of
  \(K(C)\) and the equality is the identity on the underlying rational
  functions. These section-wise identifications are induced by the canonical
  inclusion \(\struct C \hookrightarrow \mathcal K_C\), so they are compatible
  with the restriction maps (each restriction is the identity on \(K(C)\)) and
  glue to a natural isomorphism of sheaves
  \(\texttt{carrierSheaf}_0 \cong \toModuleKSheaf C\)
  (\cref{def:Scheme_toModuleKSheaf}) in
  \(\mathbf{Sh}(C, \mathbf{Mod}_{\bar k})\).
\end{proof}

\begin{lemma}

\leanok
[Short exact sequence for \(D \rightsquigarrow D + [P]\)]
  \label{lem:sheafOf_ses_single_add}
  \lean{AlgebraicGeometry.Scheme.WeilDivisor.sheafOf_ses_single_add}
  \uses{def:sheafOf, def:divisor_closed_point, def:prime_divisor,
        lem:sheafOf_mono_single_add,
        lem:cokernel_sheafOf_single_add_iso_skyscraper,
        thm:sheafOf_carrierSheaf_ses_single_add,
        lem:sheafOf_eq_carrier_of_ne_zero,
        lem:carrierSheaf_zero_iso_toModuleKSheaf,
        def:shortComplex_mathlib, def:cokernel_mathlib,
        def:skyscraperSheaf_mathlib}

  % SOURCE: [Hartshorne], IV.1, p.~296 (the inductive step of the proof
  % of Theorem~1.3, ``Next, let \(D\) be any divisor, and let \(P\) be any
  % point...''). Read from
  % references/hartshorne-algebraic-geometry.pdf, PDF page 313,
  % transcribed character-by-character from the rendered scanned page
  % image (the PDF has no text layer).
  % SOURCE QUOTE: "We consider \(P\) as a closed subscheme of \(X\). Its
  % structure sheaf is a skyscraper sheaf \(k\) sitting at the point \(P\),
  % which we denote by \(k(P)\), and its ideal sheaf is \(\mathscr L(-P)\)
  % by (II, 6.18). Therefore we have an exact sequence
  % \(0 \to \mathscr L(-P) \to \mathscr O_X \to k(P) \to 0\).
  % Tensoring with \(\mathscr L(D + P)\) we get
  % \(0 \to \mathscr L(D) \to \mathscr L(D + P) \to k(P) \to 0\).
  % (Since \(\mathscr L(D + P)\) is locally free of rank 1, tensoring by
  % it does not affect the sheaf \(k(P)\).)"
  \textit{Source: Hartshorne, IV.1, p.~296, the inductive step of the
  proof of Theorem~1.3.}
  Let \(C\) be a smooth proper geometrically irreducible curve over
  \(\bar k\), let \(D \in \mathrm{Div}(C)\) be a Weil divisor, and let \(P \in C\) be
  a closed point (equivalently, a prime divisor
  \cref{def:prime_divisor} on a smooth curve). Then the inclusion
  \(\mathcal O_C(D) \hookrightarrow \mathcal O_C(D + [P])\) of
  sub-\(\mathcal O_C\)-modules of \(\mathcal K_C\) fits into a short exact
  sequence
  \[
    0 \;\longrightarrow\; \mathcal O_C(D)
    \;\longrightarrow\; \mathcal O_C(D + [P])
    \;\longrightarrow\; k(P)
    \;\longrightarrow\; 0
  \]
  in \(\mathbf{Coh}(C)\) (equivalently, in
  \(\mathbf{Sh}(C, \mathbf{Mod}_{\bar k})\)), where \(k(P)\) denotes the
  skyscraper sheaf at \(P\) with stalk \(\bar k\) (the residue field of \(P\)
  is \(\bar k\) since \(\bar k\) is algebraically closed and \(P\) is a
  \(\bar k\)-rational closed point).
\end{lemma}

% SOURCE QUOTE PROOF: "We consider \(P\) as a closed subscheme of \(X\). Its
% structure sheaf is a skyscraper sheaf \(k\) sitting at the point \(P\),
% which we denote by \(k(P)\), and its ideal sheaf is \(\mathscr L(-P)\)
% by (II, 6.18). Therefore we have an exact sequence
% \(0 \to \mathscr L(-P) \to \mathscr O_X \to k(P) \to 0\).
% Tensoring with \(\mathscr L(D + P)\) we get
% \(0 \to \mathscr L(D) \to \mathscr L(D + P) \to k(P) \to 0\).
% (Since \(\mathscr L(D + P)\) is locally free of rank 1, tensoring by
% it does not affect the sheaf \(k(P)\).)"
\begin{proof}
  \uses{lem:sheafOf_mono_single_add,
        lem:cokernel_sheafOf_single_add_iso_skyscraper,
        thm:sheafOf_carrierSheaf_ses_single_add,
        lem:sheafOf_eq_carrier_of_ne_zero,
        lem:carrierSheaf_zero_iso_toModuleKSheaf,
        def:shortComplex_mathlib, def:cokernel_mathlib,
        def:skyscraperSheaf_mathlib}
  Hartshorne obtains the sequence by tensoring the ideal-sheaf sequence
  \(0 \to \mathcal O_C(-[P]) \to \mathcal O_C \to k(P) \to 0\) with the
  locally free rank-one sheaf \(\mathcal O_C(D+[P])\). The category
  \(\mathbf{Sh}(C, \mathbf{Mod}_{\bar k})\) has \(\bar k\)-linear sheaves
  as objects rather than \(\mathcal O_C\)-module objects, so the
  tensor \(\otimes_{\mathcal O_C}\) is not its monoidal product and that
  approach does not apply. Instead we obtain the \emph{same} short exact
  sequence from the carrier-level sequence
  \cref{thm:sheafOf_carrierSheaf_ses_single_add} and transport its first two
  terms to \(\mathcal O_C(D)\) and \(\mathcal O_C(D+[P])\).

  By \cref{thm:sheafOf_carrierSheaf_ses_single_add} the canonical mono--cokernel
  pair of the sheaf-level inclusion
  \(\texttt{carrierSheaf}_D \hookrightarrow \texttt{carrierSheaf}_{[P]+D}\) is a
  short complex (\cref{def:shortComplex_mathlib})
  \(\texttt{carrierSheaf}_D \xrightarrow{\ f\ } \texttt{carrierSheaf}_{[P]+D}
  \xrightarrow{\ \pi_f\ } \operatorname{coker} f\) that is short exact: \(f\) is a
  monomorphism by \cref{lem:sheafOf_mono_single_add}, \(\pi_f\) is an
  epimorphism as a cokernel projection (\cref{def:cokernel_mathlib}), and
  exactness at the middle term is the defining property of the cokernel of a
  monomorphism. Its third term is the skyscraper \(k(P) =
  \texttt{skyscraperSheaf}\,P.\mathrm{point}\,\bar k\)
  (\cref{def:skyscraperSheaf_mathlib}) via
  \cref{lem:cokernel_sheafOf_single_add_iso_skyscraper}.

  It remains to rewrite the first two terms as \(\mathcal O_C(D)\) and
  \(\mathcal O_C(D+[P])\). We distinguish three cases by whether either divisor
  is zero; since \(D = 0\) forces \(D + [P] = [P] \neq 0\) and \(D + [P] = 0\)
  forces \(D = -[P] \neq 0\), at most one corner occurs.
  \begin{itemize}
    \item \emph{Generic case \(D \neq 0\) and \(D + [P] \neq 0\).} Both invertible
      sheaves equal their carrier sheaves by
      \cref{lem:sheafOf_eq_carrier_of_ne_zero}, so
      \(\texttt{carrierSheaf}_D = \mathcal O_C(D)\) and
      \(\texttt{carrierSheaf}_{[P]+D} = \mathcal O_C(D+[P])\) directly.
    \item \emph{Corner \(D = 0\).} Here \(\mathcal O_C(D) = \mathcal O_C(0)\) is
      the structure sheaf \(\toModuleKSheaf C\), \emph{not}
      \(\texttt{carrierSheaf}_0\); the bridge
      \cref{lem:carrierSheaf_zero_iso_toModuleKSheaf} supplies the isomorphism
      \(\texttt{carrierSheaf}_0 \cong \toModuleKSheaf C\), identifying the first
      term with \(\mathcal O_C(0)\). The second term \(D + [P] = [P] \neq 0\) is
      rewritten by \cref{lem:sheafOf_eq_carrier_of_ne_zero}.
    \item \emph{Corner \(D + [P] = 0\)} (that is \(D = -[P]\)). Here
      \(\mathcal O_C(D+[P]) = \mathcal O_C(0) = \toModuleKSheaf C\), and the
      bridge \cref{lem:carrierSheaf_zero_iso_toModuleKSheaf} identifies the
      second term with \(\mathcal O_C(0)\). The first term \(D = -[P] \neq 0\) is
      rewritten by \cref{lem:sheafOf_eq_carrier_of_ne_zero}.
  \end{itemize}
  In every case the carrier-level short exact sequence transports to the
  asserted short exact sequence
  \[
    0 \;\longrightarrow\; \mathcal O_C(D)
    \;\longrightarrow\; \mathcal O_C(D + [P])
    \;\longrightarrow\; k(P)
    \;\longrightarrow\; 0
  \]
  in \(\mathbf{Sh}(C, \mathbf{Mod}_{\bar k})\). Concretely the first map is
  the inclusion of subsheaves of \(\mathcal K_C\) (the order constraint
  \(\ord_P(f) \geq -n_P^{D}\) implies the weaker
  \(\ord_P(f) \geq -n_P^{D} - 1\), identity on \(K(C)\)) and the second is
  the residue (``leading Laurent coefficient at the order-\((n_P^{D}+1)\)
  pole at \(P\)'') evaluation, well-defined modulo
  \(\Gamma(U, \mathcal O_C(D))\).
\end{proof}

\section{Sheaf-property correctness}
\label{sec:sheaf-property-correctness}

The construction of \cref{def:sheafOf} is presheaf-by-hand: a
per-open \(\bar k\)-submodule of \(K(C)\), with restriction along
\(V \hookrightarrow U\) given by the identity on \(K(C)\). This section
sketches the verifications that make the construction land in the
target category \(\mathbf{Sh}(C, \mathbf{Mod}_{\bar k})\).

\paragraph{Functoriality of the carrier.} For \(V \subseteq U\), the
constraint set defining \(\Gamma(U, \mathcal O_C(D))\) is a \emph{subset}
of the constraint set defining \(\Gamma(V, \mathcal O_C(D))\) (fewer
prime divisors to constrain at), so the identity on \(K(C)\) restricts
to a well-defined inclusion
\(\Gamma(U, \mathcal O_C(D)) \hookrightarrow \Gamma(V, \mathcal O_C(D))\).
Composition \(W \subseteq V \subseteq U\) is again identity-on-\(K(C)\), so
strict functoriality holds.

\paragraph{\(\bar k\)-linearity.} A \(\bar k\)-scalar \(\lambda \in \bar k\)
acts on \(K(C)\) by multiplication, and
\(\ord_Q(\lambda \cdot f) = \ord_Q(f)\) for \(\lambda \neq 0\)
(units have order zero); for \(\lambda = 0\), \(\lambda \cdot f = 0\) is
included by convention. Thus the constraint set is a
\(\bar k\)-submodule and the restriction maps are \(\bar k\)-linear.

\paragraph{Sheaf property.} The order condition
``\(\ord_Q(f) \geq -n_Q\) for every prime divisor \(Q \in U\)'' is
\emph{stalk-local} at each \(Q\): it depends only on the germ of \(f\) at
the point of \(C\) underlying \(Q\), equivalently on the image of \(f\) in
the DVR \(\mathcal O_{C, Q}\). Consequently the sheaf property reduces
to two ingredients:
\begin{enumerate}
  \item \emph{The presheaf-of-sets underlying \(\mathcal O_C(D)\) is a
    subpresheaf of the constant presheaf at \(K(C)\).} The constant
    presheaf at \(K(C)\) is a sheaf (every locally-constant section of a
    constant presheaf on a topological space is globally constant when
    the space is irreducible, which the topological space underlying an
    integral scheme is). A subpresheaf of a sheaf is itself a sheaf iff
    its membership condition is local; the order conditions are local
    at each prime divisor's underlying point, so the membership
    condition is local across an open cover and a section glued from
    locally-defined sections inherits the order condition at every
    prime divisor.
  \item \emph{The forget-functor bridge from
    \(\mathbf{Sh}(C, \mathbf{Mod}_{\bar k})\) to
    \(\mathbf{Sh}(C, \mathbf{Set})\).} The categorical forget
    \(\mathbf{Mod}_{\bar k} \to \mathbf{Set}\) creates limits, hence
    reflects the sheaf property: a presheaf of \(\bar k\)-modules is a
    sheaf iff its underlying presheaf of sets is. Combined with~(i),
    this delivers the sheaf property for the
    \(\bar k\)-module-valued \(\mathcal O_C(D)\).
\end{enumerate}

\paragraph{Invertibility.} On any open \(U \subseteq C\) disjoint from
the support of \(D\), the constraint set is exactly the structure-sheaf
constraint \(\ord_Q(f) \geq 0\) for every prime divisor \(Q \in U\); hence
\(\mathcal O_C(D)|_U = \mathcal O_C|_U\), and the trivialisation
\(\mathcal O_C(D)|_U \cong \mathcal O_C|_U\) is the identity. On a small
open \(U \subseteq C\) containing exactly one prime divisor \(Q\) from the
support of \(D\), the local generator \(\pi_Q^{-n_Q} \in K(C)\) trivialises
\(\mathcal O_C(D)|_U\) as a free rank-\(1\) \(\mathcal O_C|_U\)-module: a
section \(f \in \Gamma(U, \mathcal O_C(D))\) satisfies $\ord_Q(f) \geq
-n_Q$ iff $\pi_Q^{n_Q} \cdot f \in \Gamma(U, \mathcal O_C)$, and the
``multiplication by \(\pi_Q^{n_Q}\)'' map is an
\(\mathcal O_C|_U\)-linear isomorphism
\(\mathcal O_C(D)|_U \xrightarrow{\sim} \mathcal O_C|_U\). The two types
of trivialising opens cover \(C\) (any prime divisor lies in at most
finitely many supports across a covering, and the support of \(D\) is
finite); patching the trivialisations on overlaps amounts to changing
uniformiser, which is unit-times the original (by the DVR property), so
the cocycle is in \(\mathcal O_C^{\times}\) and the sheaf is invertible.

\section{Out of scope}

\begin{itemize}
  \item \textbf{Linear-equivalence isomorphism.} The natural
    isomorphism \(\mathcal O_C(D) \cong \mathcal O_C(D')\) when
    \(D \sim D'\) (i.e.\ \(D - D' = \div(f)\) for some
    \(f \in K(C)^{\times}\)) is \emph{not} formalised in this chapter.
    It is not needed for the main genus-zero proof: \texttt{RR.2}
    consumes only \(\chi(\mathcal O_C(D))\) and the SES additivity at
    closed points; \texttt{RR.3} consumes only the closed-point
    specialisation; \texttt{RR.4} consumes the function-field /
    degree bridge, which is divisorial in \(\mathrm{Div}(C)\) rather than
    in \(\mathrm{Pic}(C) = \mathrm{Div}(C) / \sim\). The
    Cartier-divisor \(\leftrightarrow\) \texttt{Pic} isomorphism is
    queued for the subsequent chapter.
  \item \textbf{General sheaf-multiplication
    $\mathcal O_C(D) \otimes \mathcal O_C(D') \cong \mathcal O_C(D +
    D')$.} The structure-sheaf tensor $\otimes_{\mathcal O_C}$ is not the
    monoidal product of the target category
    \(\mathbf{Sh}(C, \mathbf{Mod}_{\bar k})\) (whose objects are
    \(\bar k\)-linear sheaves rather than \(\mathcal O_C\)-module objects),
    so the inductive-step sequence is presented intrinsically by the
    cokernel route of \cref{lem:sheafOf_ses_single_add} (the subsheaf
    inclusion \cref{lem:sheafOf_mono_single_add} together with the cokernel
    identification \cref{lem:cokernel_sheafOf_single_add_iso_skyscraper})
    rather than by tensoring. The general tensor-additivity statement
    $\mathcal O_C(D) \otimes \mathcal O_C(D') \cong \mathcal O_C(D + D')$
    for arbitrary divisors \(D, D'\) is not established in this chapter;
    it is queued for the a subsequent chapter.
  \item \textbf{Pullback of \(\mathcal O_C(D)\) along a morphism.} The
    pullback \(f^* \mathcal O_C(D) \cong \mathcal O_{C'}(f^* D)\) for
    a morphism of curves \(f : C' \to C\), used in degree-computation
    arguments on a base change, is \emph{not} formalised in this
    chapter and is not needed for the main genus-zero proof.
  \item \textbf{Global-sections identification as the Riemann--Roch
    space.} The identification
    \(H^0(C, \mathcal O_C(D)) \cong L(D)\) for a general Weil divisor
    \(D\) (i.e.\ the analogue of
    \cref{lem:lineBundleAtClosedPoint_globalSections_iff} for general
    \(D\)) is queued for the a subsequent chapter. The closed-point
    case \(D = [P]\) is the only specialisation consumed by the
    main genus-zero proof, and it is handled in
    \texttt{RR.3} via \cref{lem:sheafOf_singlePoint}.
  \item \textbf{Serre duality and the canonical sheaf.} The dualizing
    sheaf \(\omega_X\) and the Serre duality theorem are not used in this
    chapter and are not required for the main genus-zero proof.
\end{itemize}
